\documentclass[11pt]{article}
\pdfinfoomitdate=1
\pdftrailerid{}
\pdfsuppressptexinfo=15
\usepackage[T1]{fontenc}
\usepackage{lmodern,amsmath,amssymb,amsthm}
\pdfglyphtounicode{parenleftbig}{0028}
\pdfglyphtounicode{parenrightbig}{0029}
\pdfgentounicode=1
\usepackage[margin=1in]{geometry}
\usepackage{microtype,enumitem,needspace}
\usepackage[numbers,sort&compress]{natbib}
\usepackage[colorlinks=true,linkcolor=blue,citecolor=blue,urlcolor=blue]{hyperref}
\newcommand{\doi}[1]{doi: \href{https://doi.org/#1}{\nolinkurl{#1}}}
\newtheorem{theorem}{Theorem}[section]
\newtheorem{lemma}[theorem]{Lemma}
\newtheorem{proposition}[theorem]{Proposition}
\newtheorem{corollary}[theorem]{Corollary}
\theoremstyle{remark}
\newtheorem{remark}[theorem]{Remark}
\newcommand{\R}{\mathcal R}
\newcommand{\Rk}{\mathcal R_k}
\newcommand{\hc}{\operatorname{hc}}
\newcommand{\girth}{\operatorname{g}}
\numberwithin{equation}{section}
\setlist{itemsep=3pt,topsep=5pt}
\hypersetup{
 pdftitle={Hamiltonian Thresholds for Planar Recolouring Graphs},
 pdfauthor={},
 pdfsubject={Hamiltonian and Hamilton-connected recolouring graphs}
}

\title{Hamiltonian Thresholds for Planar Recolouring Graphs}
\author{}
\date{October 2026}

\begin{document}
\maketitle

\begin{abstract}
Let $\Rk(G)$ be the graph of all proper colourings of a finite graph $G$ with a labelled palette $[k]$, with one-vertex recolourings as edges.  We collect sharp Hamiltonian and Hamilton-connected thresholds for planar graph classes.  A $d$-degenerate graph has a Hamilton-connected and vertex-pancyclic recolouring graph for every $k\ge d+3$.  The resulting Hamilton-connected thresholds are $7$ for all planar graphs, $5$ for triangle-free planar graphs (and hence for bipartite planar graphs), $4$ for planar graphs of girth at least six, for nonempty forests, and for the single graph $C_5$.  We give the exact Hamiltonian thresholds $4$ for planar girth-five graphs with at most seven distinct $5$-cycles, for maximum degree at most three.  For these girth-five classes, four-colour Hamilton-connectedness remains open.  The general planar girth-five four-colour problem remains open: the current bounds are $4\le h\le5$ and $4\le\hc\le5$.
\end{abstract}

\section{Introduction}

All graphs in this paper are finite and simple; they may be disconnected.  A forest is assigned infinite girth.  Colours are labelled, and a proper colouring need not use every colour. All graph classes in the threshold definitions below consist of nonempty graphs.  If $G$ is a graph, then $\Rk(G)$ has as vertices the proper maps $\sigma:V(G)\to[k]$, and two colourings are adjacent when they differ on exactly one vertex of $G$.

A graph is Hamiltonian if it has a spanning cycle, Hamilton-connected if every two distinct vertices are joined by a spanning path, and vertex-pancyclic if every vertex lies on a cycle of every length between three and the order of the graph.  For a class $\mathcal C$ of nonempty graphs define
\[
 h(\mathcal C)=\min\{q: \Rk(G)\text{ is Hamiltonian for all }G\in\mathcal C,\ k\ge q\},
\]
\[
 \hc(\mathcal C)=\min\{q: \Rk(G)\text{ is Hamilton-connected for all }G\in\mathcal C,\ k\ge q\},
\]
whenever the displayed set is nonempty.  The word threshold always has this uniform meaning; no monotonicity in $k$ is assumed.

The principal conclusions are summarized below.  A range such as $4\le\hc\le5$ records a proved four-colour lower bound and a proved five-colour upper bound, with the four-colour Hamilton-connected statement unresolved.

\begin{center}
\renewcommand{\arraystretch}{1.18}
\begin{tabular}{p{0.32\linewidth}c c p{0.34\linewidth}}
\hline
Graph class & $h$ & $\hc$ & additional conclusion\\
\hline
All planar graphs & $7$ & $7$ & $k\ge8$: vertex-pancyclic\\
Triangle-free planar graphs & $5$ & $5$ & $k\ge6$: vertex-pancyclic\\
Bipartite planar graphs & $5$ & $5$ & the cube gives the four-colour lower bound\\
Planar graphs with $\girth(G)\ge6$ & $4$ & $4$ & $k\ge5$: vertex-pancyclic\\
Nonempty forests & $4$ & $4$ & $k\ge4$: vertex-pancyclic\\
The single graph $C_5$ & $4$ & $4$ & $\R_3(C_5)$ is two $15$-cycles\\
Planar $\girth\ge5$, at most seven distinct $5$-cycles & $4$ & $4\le\hc\le5$ & four-colour Hamiltonian\\
Planar $\girth\ge5$, $\Delta\le3$ & $4$ & $4\le\hc\le5$ & four-colour Hamiltonian\\
General planar $\girth\ge5$ & $4\le h\le5$ & $4\le\hc\le5$ & the four-colour problem is open\\
\hline
\end{tabular}
\end{center}

The four-colour recolouring graph of every planar graph of girth at least five is connected by Bartier et al.~\cite{BBFHMP23}.  Earlier recolouring and Gray-code work gives useful connectivity context~\cite{DFFV06,ChMa11}, but the Hamilton-connected strengthening used here is separate.  Connectivity, a Hamilton cycle in a restricted fibre, and a finite state computation do not imply a Hamilton cycle of the complete recolouring graph.

\section{The degenerate theorem}
\label{sec:degenerate}

We begin with the general result that supplies all large-palette statements.

\begin{lemma}[Cyclic list selection]
\label{lem:list-cycle}
Let $m\ge3$ and let $L_1,\ldots,L_m$ be finite sets of size at least two.  Unless $m$ is odd and all $L_i$ are the same two-element set, there are $x_i\in L_i$ with $x_i\ne x_{i+1}$ cyclically.
\end{lemma}

\begin{proof}
If some $L_i$ has at least three elements, choose successively along the path obtained by deleting that position, then choose the remaining element to avoid both end values.  If all sets are binary and two consecutive sets differ, start with a value in their symmetric difference and continue greedily.  If all are the same binary set, alternation is possible exactly when $m$ is even. \qed
\end{proof}

\begin{lemma}[One-vertex lift]
\label{lem:single-lift}
Let $H=G-v\ne\varnothing$ and suppose $k\ge d_G(v)+3$.  If $\Rk(H)$ has a Hamilton cycle, then $\Rk(G)$ has a Hamilton cycle.
\end{lemma}

\begin{proof}
Write a Hamilton cycle of $\Rk(H)$ as $f_1,\ldots,f_m,f_1$ and put $L_i=[k]\setminus f_i(N_G(v))$.  Then $|L_i|\ge3$ and $|L_i\cap L_{i+1}|\ge2$.  Lemma~\ref{lem:list-cycle} gives cyclic choices $x_i\in L_i\cap L_{i+1}$ with consecutive values distinct.  The extensions above $f_i$ form a clique on $L_i$; traverse that clique by a Hamilton path from $x_{i-1}$ to $x_i$, and use the edge with $v$ fixed at $x_i$ to cross to the next fibre.  The exceptional all-identical binary case is impossible because every proper colouring of $H$ occurs on the base cycle and a global colour permutation changes a neighbour of $v$ to either binary colour.  The fibre paths therefore concatenate to a spanning cycle. \qed
\end{proof}


\begin{lemma}[Bridge colours]\label{lem:bridge-colours}
Let $t\ge1$, and let $S_0,\ldots,S_{t-1}$ be finite sets of colours
such that $|S_i|\ge2$ for every $i$. Suppose that $|S_j|\ge3$ for
some $j\in\{0,\ldots,t-1\}$. For any colours $a$ and $b$, there exist
colours $x_i\in S_i$ satisfying
\[
x_0\ne a,\qquad
x_i\ne x_{i-1}\quad(1\le i\le t-1),\qquad
x_{t-1}\ne b.
\]
\end{lemma}

\begin{proof}
Set $x_{-1}=a$ and $x_t=b$. Choose $x_0,\ldots,x_{j-1}$ in this
order, each time avoiding the preceding colour. Then choose
$x_{t-1},\ldots,x_{j+1}$ in this order, each time avoiding the following
colour. Each of these choices excludes at most one colour from a list of
size at least two. Finally choose
\[
x_j\in S_j\setminus\{x_{j-1},x_{j+1}\}.
\]
At most two colours are excluded, so this choice is possible. The
boundary assignments $x_{-1}=a$ and $x_t=b$ also cover $j=0$,
$j=t-1$, and $t=1$.
\end{proof}


\begin{theorem}[Degenerate Hamilton-connectedness]\label{thm:degenerate-hc}
Let $d\ge0$ be an integer, let $G$ be a nonempty $d$-degenerate graph,
and let $k\ge d+3$ be an integer. Then $\Rk(G)$ is
Hamilton-connected.
\end{theorem}

\begin{proof}
We proceed by induction on $n=|V(G)|$. If $n=1$, then
$\Rk(G)\cong K_k$, which is Hamilton-connected.
If $n=2$, the edgeless graph gives $K_k\square K_k$, which is
Hamilton-connected by Lemma~\ref{lem:product}.  For $G=K_2$, we have
$d\ge1$ and hence $k\ge4$, so the arrangement-graph result below applies.
We may therefore assume that $n\ge3$.

We next handle complete graphs $G=K_n$. The arrangement
graph $A_{k,n}$ has as vertices the ordered $n$-tuples of distinct
elements of $\{1,\ldots,k\}$, with two tuples adjacent if they differ
in exactly one coordinate. Hence
\[
\Rk(K_n)\cong A_{k,n}.
\]
Since $n-1\le d$, we have $k\ge d+3\ge n+2$.
By the Hamiltonian connectivity result of Hsu, Li, Tan and
Hsu~\cite[Theorem~1]{HLTH04}, $A_{k,n}$ is Hamilton-connected.

We may therefore assume that $G$ is not complete. Choose a vertex
$v_1$ of minimum degree in $G$, and put
\[
H:=G-v_1.
\]
The $d$-degeneracy of $G$ gives $\deg_G(v_1)\le d$. Since $G$ is
not complete, its minimum degree is at most $n-2$, so there exists
\[
w\in V(H)\setminus N_G(v_1).
\]
In particular, $H$ is nonempty. It is again $d$-degenerate, so by
the induction hypothesis $\Rk(H)$ is Hamilton-connected.

Let $f$ be a $k$-colouring of $H$. We define
\[
L(f):=\{1,\ldots,k\}\setminus f(N_G(v_1)).
\]
Thus $L(f)$ is precisely the set of colours that may be assigned to $v_1$ while extending $f$ to a proper $k$-colouring of $G$. 
Since $v_1$ has at most $d$ neighbors in $G$, we have
\[
|L(f)|\geq k-d\geq 3.
\]
For each $f\in V(\Rk(H))$, define the clique
\[
F_f:=\{f(c):c\in L(f)\}.
\]
Here $f(c)$ denotes the $k$-colouring of $G$ that restricts to $f$ on $H$ and assigns the colour $c$ to $v_1$. Thus $F_f$ is a complete graph on $|L(f)|$ vertices.

\smallskip

\noindent\textit{Observation.}
Let $f$ and $g$ be adjacent vertices of $\Rk(H)$. Then $f$ and $g$ differ on exactly one vertex $u\in V(H)$. If $u\notin N_G(v_1)$, then
\[
L(f)=L(g).
\]
If $u\in N_G(v_1)$, then the forbidden colour set for $v_1$ loses at
most one colour and gains at most one colour, and hence
\[
|L(f)\cap L(g)|\geq |L(f)|-1.
\]
In particular,
\[
|L(f)\cap L(g)|\geq 2.
\]
The edges between the cliques $F_f$ and $F_g$ are precisely the edges
\[
f(c)g(c)\qquad \text{with } c\in L(f)\cap L(g).
\]

\smallskip

\noindent\textit{Lifting claim.}
Let
\[
P:f_0,f_1,\dots,f_t
\]
be a path in $\Rk(H)$ with $t\ge1$, and put
\[
S_i:=L(f_i)\cap L(f_{i+1})\qquad(0\le i\le t-1).
\]
Suppose that $|S_j|\ge3$ for at least one $j$. Given colours
\[
a\in L(f_0),\qquad b\in L(f_t),
\]
there is a Hamilton path in
\[
\Rk(G)\bigg[\bigcup_{i=0}^t F_{f_i}\bigg]
\]
from $f_0(a)$ to $f_t(b)$.

Indeed, by the observation, $|S_i|\ge2$ for every $i$.
Lemma~\ref{lem:bridge-colours} therefore provides bridge colours
$x_i\in S_i$ such that
\[
x_0\neq a,\qquad x_i\neq x_{i-1}\ \ (1\leq i\leq t-1),\qquad x_{t-1}\neq b.
\]
Since each $F_{f_i}$ is complete, choose Hamilton paths
\begin{align*}
P_0&:\ f_0(a)\leadsto f_0(x_0)\quad \text{in }F_{f_0},\\
P_i&:\ f_i(x_{i-1})\leadsto f_i(x_i)\quad \text{in }F_{f_i}\qquad (1\leq i\leq t-1),\\
P_t&:\ f_t(x_{t-1})\leadsto f_t(b)\quad \text{in }F_{f_t}.
\end{align*}
For each $i\in\{0,\dots,t-1\}$, the vertices $f_i(x_i)$ and $f_{i+1}(x_i)$ are adjacent in $\Rk(G)$. Concatenating the paths $P_0,P_1,\dots,P_t$ along these bridge edges gives the required Hamilton path.

We now verify the additional hypothesis for the paths used below.
Since $H$ is $k$-colourable, permuting the colour names in a proper
colouring shows that $w$ takes all $k$ colours among the vertices of
$\Rk(H)$. In particular, $|V(\Rk(H))|\ge k\ge3$.
After deleting at most one colouring, at least $k-1\ge2$ different
colours still occur at $w$. Thus every path containing all but at most
one vertex of $\Rk(H)$ has a step $f_jf_{j+1}$ that changes
the colour of $w$. Since $w\notin N_G(v_1)$, this step satisfies
\[
L(f_j)=L(f_{j+1}),\qquad
|S_j|=|L(f_j)|\ge3.
\]
Consequently, the lifting claim applies to every such path.

Let
\[
x=\alpha(a)\qquad \text{and} \qquad y=\beta(b)
\]
be two distinct vertices of $\Rk(G)$, where $\alpha,\beta\in V(\Rk(H))$, $a\in L(\alpha)$, and $b\in L(\beta)$. 
We will show that there is a Hamilton path from $x$ to $y$ in $\Rk(G)$.

\smallskip

\noindent\textit{Case 1: $\alpha\neq \beta$.}
Since $\Rk(H)$ is Hamilton-connected, there is a Hamilton path
\[
\alpha=f_0,f_1,\dots,f_t=\beta
\]
in $\Rk(H)$. Since it contains every vertex of
$\Rk(H)$, it satisfies the additional hypothesis of the
lifting claim.

Applying the lifting claim to the path
\[
f_0,f_1,\dots,f_t
\]
with the endpoint colours $a\in L(f_0)$ and $b\in L(f_t)$, we obtain a Hamilton path in
\[
\Rk(G)\bigg[\bigcup_{i=0}^{t}F_{f_i}\bigg]
\]
from $f_0(a)=x$ to $f_t(b)=y$. Since $f_0,\dots,f_t$ exhaust all vertices of $\Rk(H)$, the cliques
\[
F_{f_0},F_{f_1},\dots,F_{f_t}
\]
partition $V(\Rk(G))$. Therefore this path is a Hamilton path in $\Rk(G)$ from $x$ to $y$.

\smallskip

\noindent\textit{Case 2: $\alpha=\beta$.}
Since $x\neq y$, we have $a\neq b$. Write
\[
f^\ast:=\alpha,\qquad c_1:=a,\qquad c_2:=b.
\]
Suppose, for a contradiction, that there is no Hamilton path in $\Rk(G)$ from
\[
x=f^\ast(c_1)\qquad \text{to}\qquad y=f^\ast(c_2).
\]

We first show that $F_{f^\ast}$ must have size exactly three. We then
prove that one of its three vertices, namely $f^\ast(c_3)$, has no
neighbor outside $F_{f^\ast}$. Recolouring a vertex of degree at most
$d$ in $H$ will give a contradiction.

\smallskip

\noindent\textit{Claim 1.} We have $|F_{f^\ast}|=3$.

Suppose instead that $|F_{f^\ast}|\geq 4$. Since $\Rk(H)$ is Hamilton-connected and has at least three vertices, every edge of $\Rk(H)$ lies on a Hamilton cycle: indeed, if $fg$ is an edge, then a Hamilton path from $f$ to $g$ together with the edge $fg$ forms a Hamilton cycle. Choose a Hamilton cycle $C$ in $\Rk(H)$ containing $f^\ast$, and let $f_i$ and $f_j$ be the two neighbors of $f^\ast$ on $C$. Then $C-f^\ast$ is a Hamilton path from $f_i$ to $f_j$ in $\Rk(H)-f^\ast$.

By the observation,
\[
|L(f^\ast)\cap L(f_i)|\geq |L(f^\ast)|-1=|F_{f^\ast}|-1\geq 3,
\]
and similarly
\[
|L(f^\ast)\cap L(f_j)|\geq 3.
\]
Choose
\[
p\in L(f^\ast)\cap L(f_i)\setminus\{c_1,c_2\},
\]
and then choose
\[
q\in L(f^\ast)\cap L(f_j)\setminus\{c_1,p\}.
\]
Then
\[
p\neq q,\qquad p\notin\{c_1,c_2\},\qquad q\neq c_1.
\]

Since $F_{f^\ast}$ is a complete graph on at least four vertices, we can partition its vertex set into two vertex-disjoint paths, one with ends $f^\ast(c_1)$ and $f^\ast(p)$, and the other with ends $f^\ast(q)$ and $f^\ast(c_2)$. Denote these two paths by $P$ and $P'$ respectively. If $q=c_2$, take $P'$ to be the single vertex $f^\ast(c_2)$.

The path $C-f^\ast$ contains all but one vertex of
$\Rk(H)$, so it satisfies the additional hypothesis of the
lifting claim. Apply that claim with endpoint colours $p$ and $q$.
This yields a Hamilton path
\[
Q
\]
in
\[
\Rk(G)\bigg[\bigcup_{f\neq f^\ast}F_f\bigg]
\]
from $f_i(p)$ to $f_j(q)$. Using the bridge edges
\[
f^\ast(p)f_i(p)\qquad\text{and}\qquad f_j(q)f^\ast(q),
\]
the concatenation
\[
P,\ Q,\ P'
\]
is a Hamilton path in $\Rk(G)$ from $x$ to $y$, a contradiction. Hence $|F_{f^\ast}|=3$.

We may therefore write
\[
L(f^\ast)=\{c_1,c_2,c_3\}.
\]
Since $|L(f^\ast)|\geq k-d$, this implies $k=d+3$.

\smallskip

\noindent\textit{Claim 2.} The vertex $f^\ast(c_3)$ has no neighbor outside $F_{f^\ast}$.

Suppose not. Then there is a colouring $f_i$ adjacent to $f^\ast$ in $\Rk(H)$ such that
\[
c_3\in L(f_i).
\]
Choose a Hamilton cycle $C$ of $\Rk(H)$ containing the edge $f^\ast f_i$, and let $f_j$ be the other neighbor of $f^\ast$ on $C$. Then $C-f^\ast$ is a Hamilton path from $f_i$ to $f_j$ in $\Rk(H)-f^\ast$.

By the observation, $|L(f^\ast)\cap L(f_j)|\geq 2$. Since $L(f^\ast)=\{c_1,c_2,c_3\}$, choose
\[
q\in L(f^\ast)\cap L(f_j)\setminus\{c_3\}.
\]
Thus $q\in\{c_1,c_2\}$.

Again, $C-f^\ast$ contains all but one vertex of $\Rk(H)$.
Applying the lifting claim with endpoint colours $c_3$ and $q$, we
obtain a Hamilton path
\[
Q
\]
through all cliques other than $F_{f^\ast}$, from $f_i(c_3)$ to $f_j(q)$.

If $q=c_2$, then
\[
f^\ast(c_1),\ f^\ast(c_3),\ Q,\ f^\ast(c_2)
\]
using the bridge edges $f^\ast(c_3)f_i(c_3)$ and $f_j(c_2)f^\ast(c_2)$ is a Hamilton path from $x$ to $y$, a contradiction. If $q=c_1$, then the same construction gives a Hamilton path from $y$ to $x$, and reversing it yields a Hamilton path from $x$ to $y$, again a contradiction. Hence $f^\ast(c_3)$ has no neighbor outside $F_{f^\ast}$.

Since $H$ is nonempty and $d$-degenerate, choose $u\in V(H)$ with
$\deg_H(u)\le d$. Under $f^\ast$, at least $k-d\ge3$ colours are
absent from $N_H(u)$. We can therefore choose one of these colours,
say $\gamma$, different from both $f^\ast(u)$ and $c_3$.
Recolouring $u$ with $\gamma$ gives a proper colouring $g$ of $H$
adjacent to $f^\ast$. No neighbor of $v_1$ has acquired the colour
$c_3$, so $c_3\in L(g)$. Thus $f^\ast(c_3)$ is adjacent to
$g(c_3)$ outside $F_{f^\ast}$, contradicting Claim 2.

In both cases there is a Hamilton path from $x$ to $y$. Therefore $\Rk(G)$ is Hamilton-connected.
\end{proof}

\begin{theorem}[Degenerate vertex-pancyclicity]\label{thm:degenerate-vp}
Let $d\ge0$ be an integer, let $G$ be a nonempty $d$-degenerate graph,
and let $k\ge d+3$ be an integer.
Then the recolouring graph $\Rk(G)$ is vertex-pancyclic.

\end{theorem}

\begin{proof}
We proceed by induction on $|V(G)|$.

If $|V(G)|=1$, then $\Rk(G)\cong K_k$. Since $k\ge 3$,
every vertex of $K_k$ lies on a cycle of every length from $3$ to $k$.
Thus the result holds.

Assume now that $|V(G)|\ge 2$, and suppose the result holds for all
$d$-degenerate graphs with fewer vertices.

Choose a vertex $v_1$ of degree at most $d$, and put
\[
H:=G-v_1.
\]
For each proper $k$-colouring $f$ of $H$, define
\[
L(f):=\{1,\ldots,k\}\setminus f(N_G(v_1)).
\]
Thus $L(f)$ is the set of colours that can be assigned to $v_1$ while
extending $f$ to a proper $k$-colouring of $G$. Since
$|N_G(v_1)|\le d$, we have
\[
|L(f)|\ge k-d\ge 3.
\]
For each $f\in V(\Rk(H))$, define
\[
F_f:=\{f(c):c\in L(f)\}.
\]
Then $F_f$ is a clique in $\Rk(G)$.

Let
\[
X:=\Rk(G),\qquad Y:=\Rk(H).
\]
The sets $F_f$, where $f\in V(Y)$, form a partition of $V(X)$.

As observed in the preceding proof, if $f$ and $g$ are adjacent in $Y$,
then $|L(f)\cap L(g)|\ge 2$, and the edges between $F_f$ and $F_g$ are
precisely $f(c)g(c)$ with $c\in L(f)\cap L(g)$.

Now fix an arbitrary vertex of $X$,
\[
x=\alpha(a),
\]
where $\alpha\in V(Y)$ and $a\in L(\alpha)$.
We prove that $x$ lies on a cycle of every possible length.

Let
\[
M:=|V(Y)|,\qquad N:=|V(X)|.
\]

\smallskip

\noindent\textit{Closed lifting claim.}
Let
\[
C:f_0,f_1,\ldots,f_{s-1},f_0
\]
be a cycle in $Y$ with $s\ge3$ and $f_0=\alpha$. Then, for every integer $\ell$
satisfying
\[
2s+1\le \ell\le \sum_{i=0}^{s-1}|F_{f_i}|,
\]
there is a cycle in $X$ through $x$ of length $\ell$ using only vertices
from the cliques $F_{f_0},F_{f_1},\ldots,F_{f_{s-1}}$.

For $i$ modulo $s$, put
\[
S_i:=L(f_i)\cap L(f_{i+1}).
\]
We first show that there are bridge colours $y_i\in S_i$ such that
$y_{i-1}\ne y_i$ for every $i$ modulo $s$.

We use an elementary fact about lists on a cycle. When all lists have
size at least two, a proper choice of colours exists unless the cycle
is odd and all lists are the same two-element set. To see this, if
some list has size at least three, leave it until last and greedily
colour the path formed by the remaining positions. The last position
has at most two forbidden colours. If all lists have size two but are
not identical, choose adjacent positions whose lists differ, say
$S_r$ and $S_{r-1}$, and start with a colour in
$S_r\setminus S_{r-1}$. Greedily colour the positions
$r+1,\ldots,r-1$ in cyclic order. The last colour automatically
differs from the first, since the first colour is absent from
$S_{r-1}$. Finally, if all lists are the same two-element set and the
cycle is even, alternate the two colours.

It remains to exclude the exceptional odd-cycle configuration in our
setting. Suppose that $s$ is odd and
\[
S_0=S_1=\cdots=S_{s-1}=\{p,q\}.
\]
For every $i$, the observation gives
$2=|S_i|\ge |L(f_i)|-1$, so $|L(f_i)|=3$. Consequently,
\[
k-3=|f_i(N_G(v_1))|\le \deg_G(v_1)\le d\le k-3.
\]
All inequalities are equalities. Thus $k=d+3$, the vertex $v_1$ has
exactly $d$ neighbors, and these neighbors receive pairwise distinct
colours under every $f_i$. None of these colours is $p$ or $q$.

Every step $f_if_{i+1}$ changes a vertex in $N_G(v_1)$; otherwise
$L(f_i)=L(f_{i+1})$ would imply $|S_i|=3$. In particular, $d\ge1$.
Fix an ordering $u_1,\ldots,u_d$ of $N_G(v_1)$, and let
\[
A:=\{1,\ldots,k\}\setminus\{p,q\}.
\]
For each $i$, exactly one colour $r_i\in A$ is absent from
$f_i(N_G(v_1))$. Hence
\[
\bigl(f_i(u_1),\ldots,f_i(u_d),r_i\bigr)
\]
is a permutation of the $d+1$ colours in $A$. At the step
$f_if_{i+1}$, the changed neighbor must receive $r_i$, because the
other neighbor colours remain fixed and pairwise distinct. Its old
colour becomes $r_{i+1}$. Thus this step exchanges the entries at one
neighbor position and the last position, reversing the parity of the
permutation. Since $f_s=f_0$, an odd number of such exchanges cannot
return to the initial permutation. This contradiction excludes the
exceptional configuration, and proves the existence of the colours
$y_i$.

The integer intervals
\[
[3,|F_{f_0}|],\ [2,|F_{f_1}|],\ldots,[2,|F_{f_{s-1}}|]
\]
have sum equal to the full integer interval
$[2s+1,\sum_i |F_{f_i}|]$. We may therefore choose integers
$m_0,\ldots,m_{s-1}$ with
\[
3\le m_0\le |F_{f_0}|,\qquad
2\le m_i\le |F_{f_i}|\quad (1\le i\le s-1),\qquad
\sum_{i=0}^{s-1}m_i=\ell.
\]
Inside $F_{f_0}$ take a path from $f_0(y_{s-1})$ to $f_0(y_0)$ on
$m_0$ vertices and containing $x$. For $1\le i\le s-1$, take inside
$F_{f_i}$ a path from $f_i(y_{i-1})$ to $f_i(y_i)$ on $m_i$ vertices.
The inequalities $y_{i-1}\ne y_i$ make these paths possible, and the
bridge edges $f_i(y_i)f_{i+1}(y_i)$, with indices modulo $s$, close
them into the desired cycle.

\smallskip

Since $|L(\alpha)|\ge3$, the clique $F_\alpha$ contains a triangle
through $x$.

Next we show that $x$ lies on cycles of lengths $4,5,6$.
Since $H$ is nonempty and $d$-degenerate, choose a vertex $u\in V(H)$
of degree at most $d$ in $H$. Because $k\ge d+3$, we can recolour $u$
to some colour different from its present colour, avoiding all colours on
$N_H(u)$ and, if $u\in N_G(v_1)$, also avoiding the colour $a$.
Thus there is a colouring $\beta$ adjacent to $\alpha$ in $Y$ such that
\[
a\in L(\beta).
\]
Since
\[
|L(\alpha)\cap L(\beta)|\ge 2,
\]
choose
\[
p\in L(\alpha)\cap L(\beta),\qquad p\ne a.
\]
Also choose
\[
q\in L(\alpha)\setminus\{a,p\}
\]
and
\[
r\in L(\beta)\setminus\{a,p\}.
\]
Then the following are cycles in $X$ through $x$:
\[
\alpha(a),\alpha(p),\beta(p),\beta(a),\alpha(a),
\]
of length $4$;
\[
\alpha(a),\alpha(q),\alpha(p),\beta(p),\beta(a),\alpha(a),
\]
of length $5$; and
\[
\alpha(a),\alpha(q),\alpha(p),\beta(p),\beta(r),\beta(a),\alpha(a),
\]
of length $6$.

It remains to construct cycles of all lengths $\ell\ge 7$.

By the induction hypothesis, $Y=\Rk(H)$ is vertex-pancyclic.
Hence, for every integer $s$ with
\[
3\le s\le M,
\]
there is a cycle of length $s$ in $Y$ containing $\alpha$.

Let
\[
C:f_0,f_1,\ldots,f_{s-1},f_0
\]
be such a cycle, where $f_0=\alpha$.
By the closed lifting argument, and since every $|F_{f_i}|\ge 3$, this
base cycle gives cycles through $x$ of every length
\[
2s+1\le \ell\le 3s.
\]
The intervals
\[
[2s+1,3s]
\]
overlap consecutively as $s$ ranges from $3$ to $M$. Hence we obtain
cycles through $x$ of every length
\[
7\le \ell\le 3M.
\]

Finally, take a Hamilton cycle
\[
C:f_0,f_1,\ldots,f_{M-1},f_0
\]
of $Y$ containing $\alpha=f_0$.
Such a cycle exists by the induction hypothesis that $Y$ is
vertex-pancyclic.
The closed lifting argument applied to this cycle gives cycles through
$x$ of every length in the interval
\[
2M+1\le \ell\le \sum_{f\in V(Y)}|F_f|=N.
\]
This interval overlaps $[7,3M]$, since $2M+1\le 3M$.
Combining all cases, $x$ lies on a cycle of every length
\[
3\le \ell\le N.
\]

Since $x$ was arbitrary, $\Rk(G)$ is vertex-pancyclic.
\end{proof}


\begin{corollary}
\label{cor:degenerate-hamiltonian}
Every nonempty $d$-degenerate graph has a Hamiltonian $\Rk(G)$ for $k\ge d+3$; it is Hamilton-connected and vertex-pancyclic in the same range.
\end{corollary}


\section{Endpoint lifting and restoration}
\label{sec:restoration}

A colouring of a deleted graph determines the available colours at the
vertices to be restored.  The resulting extension graph is a fibre.  In
all the constructions below, assignments retain their actual coordinate
values: equal assignments in two different fibres are distinct vertices,
and a cross edge joins their two copies.  In particular, a move of one
vertex of the deleted graph may change several available lists at once.

\begin{lemma}[Products of complete graphs]
\label{lem:product}
Let $s\ge1$ and $F=\square_{j=1}^s K_{A_j}$, where every $A_j$ is
finite and $|A_j|\ge2$.  If all lists have size two, $F$ has a
Hamilton path between two prescribed vertices if and only if those
vertices lie in opposite bipartition classes.  If at least one list has
size at least three, $F$ is Hamilton-connected.
\end{lemma}

\begin{proof}
In the binary case $F$ is the hypercube $Q_s$, whose bipartition
classes have equal size.  The endpoint condition is therefore necessary.
For sufficiency, induct on $s$, starting with the single edge $Q_1$.
Choose a coordinate on which the prescribed endpoints differ and split
along it.  Their projections into $Q_{s-1}$ have the same parity.  Choose
a projected vertex of the opposite parity; use the two induction paths
to and from its copies, joined by the edge between those copies.

For the second assertion, we first show that $Y\square K_r$ is
Hamilton-connected whenever $Y$ is Hamilton-connected and
$|V(Y)|\ge3$.  Induct on $r$, starting with $r=1$.  If the two
endpoints lie in different copies of $Y$, order the copies starting and
ending at those copies.  Choose successive bridge vertices in $V(Y)$,
with the first different from the first endpoint projection, consecutive
bridges different, and the last different from the last endpoint
projection.  Such choices exist because $|V(Y)|\ge3$: choose greedily,
leaving the last bridge until the end, when it has at most two forbidden
values.  Prescribed-end Hamilton paths inside the copies then join into
a spanning path.  If both endpoints lie in one copy, take a Hamilton
path between them there and choose an edge $xy$ of it.  By induction the
other $r-1$ copies have a spanning path from a copy of $x$ to a copy of
$y$ in the same one of those copies.  Insert that path in place of
$xy$, using the two cross edges.  Start with a complete-graph factor of
order at least three and add the remaining factors to obtain the claim.
\end{proof}

\begin{lemma}[Endpoint product-fibre lift]
\label{lem:path-fibre}
Let $m\ge3$, $s\ge1$, and let
\[
 F_i=\square_{j=1}^s K_{A_{i,j}},\qquad |A_{i,j}|\ge2,
 \qquad 0\le i<m,
\]
be disjoint fibres arranged along a path.  Join equal assignments in
consecutive fibres.  At every transition assume that at least one list
is unchanged and that, for every coordinate,
\[
 |A_{i,j}\setminus A_{i+1,j}|\le1,
 \qquad |A_{i+1,j}\setminus A_{i,j}|\le1.
\]
If an internal fibre $F_r$, $1\le r\le m-2$, has a factor of order
at least three, then every prescribed $p\in F_0$ and $q\in F_{m-1}$
are joined by a path covering all vertices of all fibres exactly once.
\end{lemma}

\begin{proof}
Write
\[
 J_i=\square_{j=1}^s K_{A_{i,j}\cap A_{i+1,j}}
\]
for the graph of common assignments.  The intersections are nonempty,
and an unchanged coordinate supplies an edge of $J_i$.

Starting at $p$, record the reachable exit assignments at each
interface.  If a fibre has a reachable entrance $a$, an interface edge
has an admissible Hamilton-path endpoint: in a binary fibre choose its
endpoint of parity opposite to $a$, and in a nonbinary fibre choose an
endpoint different from $a$.  Lemma~\ref{lem:product} therefore shows
that the reachable set never becomes empty.

At $F_r$, its outgoing interface $J_r$ still contains an edge after any
one assignment is removed.  Indeed, if its large coordinate is
unchanged, $J_r$ contains a triangle.  Otherwise that coordinate has
intersection size at least two, and a different unchanged coordinate
supplies a $K_2\square K_2$.  Take a reachable entrance $a$ in $F_r$.
Hamilton-connectedness of $F_r$ makes every assignment in
$V(J_r)\setminus\{a\}$ reachable; this set contains an entire edge.

An incoming reachable edge permits every outgoing assignment to be
reached in the next fibre.  In a binary fibre the two endpoints of the
edge have opposite parities, so one has the parity needed for any
specified outgoing assignment.  In a nonbinary fibre they are distinct,
so one differs from that assignment.  Thus the reachable exits contain
an edge at every subsequent interface.  At the last fibre, one endpoint
of the incoming edge can be joined to the specified $q$.  Backtracking
the finite choices supplies one Hamilton path in each fibre; the actual
cross edges concatenate them into the required spanning path.
\end{proof}

\begin{lemma}[Hamilton-connected independent-set restoration]
\label{lem:independent-hc}
Let $S$ be a nonempty independent set, $H=G-S\ne\varnothing$, and
$k\ge4$.  Suppose
\[
 d_G(u)\le k-2\quad(u\in S),\qquad
 |N_G(x)\cap S|\le |S|-1\quad(x\in V(H)).
\]
Assume that $\Rk(H)$ is Hamilton-connected and that some proper
$k$-colouring of $H$ leaves at least three colours available at a vertex
of $S$.  If $\alpha,\beta$ are proper $k$-colourings of $G$ with
$\alpha|_H\ne\beta|_H$, then $\Rk(G)$ has an
$\alpha$--$\beta$ Hamilton path.
\end{lemma}

\begin{proof}
Take a Hamilton path of $\Rk(H)$ between the two restriction
colourings.  For a base colouring $f$, the fibre is exactly
\[
 \square_{u\in S}K_{A_f(u)},\qquad
 A_f(u)=[k]\setminus f(N_G(u)).
\]
Every list is at least binary.  A move at an actual vertex $x$ of $H$
leaves the list of at least one non-neighbour in $S$ unchanged; every
other list loses and gains at most one colour.  These are the transition
conditions of Lemma~\ref{lem:path-fibre}, even when several lists
change together.

A colouring with a nonbinary fibre has at least $k$ distinct colour
permutations: fix a vertex of the nonempty $H$ and send its colour to
each of the $k$ labels.  Every resulting fibre is nonbinary.  The base
Hamilton path has only two endpoints, so one of these colourings is
internal.  Lemma~\ref{lem:path-fibre} now lifts the path with the
specified, actual colours of $\alpha$ and $\beta$ on $S$.
\end{proof}

For use of this lemma, first choose a vertex $w$ on which two prescribed
colourings differ, and then choose $S$ with $w\notin S$.  This ensures
different restriction colourings; the lemma does not claim to handle
arbitrary endpoints in the same fibre.

\begin{lemma}[Cyclic product-fibre lift]
\label{lem:fibre-lift}
Let $m\ge3$ disjoint fibres
$F_i=\square_{j=1}^sK_{A_{i,j}}$, with $s\ge1$ and
$|A_{i,j}|\ge2$, be cyclically ordered.  Join equal assignments in
consecutive fibres and define
$J_i=\square_j K_{A_{i,j}\cap A_{i+1,j}}$, with indices modulo $m$.
Suppose every $J_i$ contains an edge, $F_0$ has a factor of order at
least three, and there are distinct edges
$e_0\in E(J_0)$ and $e_{m-1}\in E(J_{m-1})$ (compared as actual
assignment pairs).  Then the layered graph has a Hamilton cycle.
\end{lemma}

\begin{proof}
Choose an edge $e_i$ in every remaining interface, and choose
$z_0\in V(e_0)\setminus V(e_{m-1})$.  For $1\le i<m$, there is a
bijection $\phi_i:V(e_{i-1})\to V(e_i)$ such that each
$z,\phi_i(z)$ is an admissible pair of Hamilton-path endpoints in
$F_i$.  In a binary fibre match opposite parities.  Otherwise use a
bijection with no fixed assignment, which always exists between two
two-element sets: interchange equal sets, send a common member to the
other target member when there is exactly one common member, and use
either bijection when the sets are disjoint.

Put $z_i=\phi_i(z_{i-1})$.  Traverse $F_i$ from $z_{i-1}$ to $z_i$
for $1\le i<m$.  Since $z_{m-1}\in V(e_{m-1})$,
$z_{m-1}\ne z_0$, so the nonbinary $F_0$ has a Hamilton path from
$z_{m-1}$ to $z_0$.  Join the fibre paths by the cross edges at the
assignments $z_i$.  They cover every fibre once and form one cycle.
\end{proof}

\begin{lemma}[Independent-set restoration]
\label{lem:independent-restoration}
Let $S$ be a nonempty independent set, $H=G-S\ne\varnothing$, and
$k\ge4$.  Suppose
\[
 d_G(u)\le k-2\quad(u\in S),\qquad
 |N_G(x)\cap S|\le |S|-1\quad(x\in V(H)).
\]
If $\Rk(H)$ is Hamiltonian and some proper $k$-colouring of $H$ leaves
at least three colours available at a vertex of $S$, then $\Rk(G)$ is
Hamiltonian.
\end{lemma}

\begin{proof}
Take a Hamilton cycle of the base and use the same product fibres as in
Lemma~\ref{lem:independent-hc}.  Every interface contains an edge.
Index a nonbinary fibre as $F_0$ and choose a large coordinate in it.
At either adjacent interface, that coordinate supplies a triangle if
unchanged, or a $K_2\square K_2$ together with a different unchanged
coordinate if it changes.  Thus both adjacent interfaces have at least
two edges, and distinct edges $e_0,e_{m-1}$ can be chosen.  Apply
Lemma~\ref{lem:fibre-lift}.
\end{proof}

\begin{lemma}[Single-vertex Hamilton-connected lift]
\label{lem:single-hc}
Let $H=G-v\ne\varnothing$, $k\ge4$, and $k\ge d_G(v)+3$.
Assume that $v$ has a non-neighbour in $G$, $\Rk(H)$ is
Hamilton-connected, and $\delta(\Rk(H))>d_G(v)$.  Then $\Rk(G)$ is
Hamilton-connected.
\end{lemma}

\begin{proof}
For a base colouring $f$, write $A_f=[k]\setminus f(N_G(v))$.
The fibre is the clique $K_{A_f}$, $|A_f|\ge3$, and adjacent base
fibres have intersections of size at least two.

Let $w$ be a non-neighbour of $v$.  Any base path covering all states,
or all but one state, has a step changing $w$.  Indeed, global colour
permutations show that $w$ takes all $k$ colours among the base states;
omitting one state leaves at least $k-1\ge3$ colours.  At a step
changing $w$, the available list for $v$ is unchanged, so its
intersection has size at least three.  Consequently such a base path
can be lifted with arbitrary prescribed first and last fibre colours.
To see this, choose bridge colours in the successive intersections,
requiring consecutive bridge colours to differ and avoiding the
prescribed endpoint colour at either end.  All intersection lists have
size at least two and one has size at least three.  Greedily choose from
the two ends towards a three-element list, and choose at that list last
while avoiding its at most two forbidden colours.  Then traverse each
clique between its distinct entrance and exit colours.

If the prescribed states lie in different fibres, lift the base
Hamilton path between their restrictions.  Suppose instead that both
lie over $f$, with distinct colours $a,b$.  A Hamilton-connected graph
with at least three vertices has a Hamilton cycle through every edge:
close a Hamilton path between the ends of the edge by that edge.
Here the base has at least $k$ states.

If $|A_f|\ge4$, choose a base Hamilton cycle through $f$, whose
neighbours on the cycle are $g,h$.  Both intersections with $A_f$
have size at least three.  Choose
$p\in A_f\cap A_g\setminus\{a,b\}$ and
$q\in A_f\cap A_h\setminus\{a,p\}$.  Split the root clique into
disjoint paths $a$--$p$ and $q$--$b$ covering it; the second may be a
single vertex if $q=b$.  The base cycle minus $f$ is a path covering
all other states, and the preceding lifting observation traverses those
fibres with endpoint colours $p,q$.  The two cross edges and the two
root paths give the desired spanning path.

If $A_f=\{a,b,c\}$, at most $d_G(v)$ base moves from $f$ make $c$
unavailable: for each neighbour of $v$, only its move to colour $c$
can do so.  The state-degree assumption gives a base edge $fg$
preserving $c$.  Choose a base Hamilton cycle containing $fg$ and
write $h$ for the other neighbour of $f$ on it.  Select
$q\in (A_f\cap A_h)\setminus\{c\}$, so $q=a$ or $b$.  If $q=b$,
start with the root edge $a c$, cross at colour $c$, lift the path
through all other base states, and return to the root at colour $b$.
If $q=a$, carry out the same construction with the two endpoints
interchanged and reverse the resulting path.  This covers the three
root states and all other fibres exactly once.
\end{proof}

A Hamilton-connected graph of order at least four has minimum degree at
least three.  A vertex of degree two would, with its neighbours chosen
as the two endpoints, force a spanning path to have only three vertices;
degrees zero and one are excluded as well.  Every nonempty colourable
$H$ has at least $k$ states in $\Rk(H)$, by colour permutation.
These observations will supply the state-degree condition in the
low-degree inductions.

\section{Hamilton-connected planar classes}
\label{sec:hc-planar}

\subsection{All planar graphs}

We first establish the marked structural statement needed to retain a
coordinate on which the prescribed endpoint colourings differ.

\begin{lemma}[Separated degree-five vertices with a marked vertex]
\label{lem:marked-five}
If $G$ is planar and $\delta(G)\ge5$, then, for any $w\in V(G)$,
there are two degree-five vertices different from $w$ whose distance is
at least three.  Vertices in different components have infinite distance.
\end{lemma}

\begin{proof}
Extend $G$, without adding vertices, to a simple plane triangulation
$T$.  A degree-five vertex in $T$ still has degree five in $G$, and a
pair at distance at least three in $T$ has that property in $G$.
Moreover $\delta(T)\ge5$.  It suffices to prove the assertion in $T$.
Euler's formula gives at least twelve degree-five vertices.

For a degree-five vertex $x$, put $W=N_T[x]$, $U=V(T)\setminus W$,
$r=|U|$, and let $B$ be the vertices at distance exactly two from $x$;
write $b=|B|$.  The five neighbours of $x$ form a cycle in their cyclic
order, so $e_T(W)\ge10$.  Contract the connected $W$ and simplify
parallel edges.  The resulting plane graph has $r+1$ vertices and
$e_T(U)+b$ edges.  As $|V(T)|\ge12$, the planar edge bound gives
$e_T(U)+b\le3r-3$.  Since $e(T)=3r+12$,
\[
 \sum_{z\in U}d_T(z)
 =e_T(U)+e(T)-e_T(W)\le6r-b-1.
\]
Let $F_x$ be the degree-five vertices at distance at least three from
$x$.  The same sum is exactly
\[
 6r-b-|F_x|
 +\sum_{z\in B}(d_T(z)-5)
 +\sum_{\substack{z\in U\setminus B\\ d_T(z)\ge6}}(d_T(z)-6).
\]
It follows that $|F_x|\ge1$, and if $|F_x|=1$, every vertex of $B$
has degree five.

Suppose every pair of degree-five vertices at distance at least three
contains $w$, and let $L$ be all degree-five vertices except $w$.
For every $x\in L$, the nonempty $F_x$ must be $\{w\}$.  Hence
$w$ is at distance at least three from $L$, and every vertex at distance
exactly two from $x$ belongs to $L$.  Take a shortest path from $w$
to $L$, ending at $x$.  The vertex two steps before $x$ is also in
$L$ and is closer to $w$, a contradiction.
\end{proof}

\begin{theorem}[Planar Hamilton-connectedness]
\label{thm:planar-hc}
For every nonempty planar graph $G$ and every $k\ge7$, $\Rk(G)$ is
Hamilton-connected.  Consequently
$h(\textup{Planar})=\hc(\textup{Planar})=7$.
\end{theorem}

\begin{proof}
Prove the seven-colour assertion by induction on $|V(G)|$, allowing
disconnected graphs throughout.  If $|V(G)|\le5$, the graph is
four-degenerate and Theorem~\ref{thm:degenerate-hc} applies.
For $|V(G)|\ge6$, a vertex $v$ of degree at most four has a
non-neighbour.  The smaller $H=G-v$ is Hamilton-connected by induction.
For every seven-colouring $f$ of $H$, which has at least five vertices,
\[
 d_{\R_7(H)}(f)
 =\sum_{x\in V(H)}(6-|f(N_H(x))|)
 \ge6|V(H)|-2|E(H)|\ge12>d_G(v).
\]
Lemma~\ref{lem:single-hc} restores $v$.

Otherwise $\delta(G)\ge5$.  For prescribed distinct colourings
$\alpha,\beta$, choose an actual vertex $w$ with
$\alpha(w)\ne\beta(w)$.  Lemma~\ref{lem:marked-five} supplies
$S=\{u,v\}$, avoiding $w$, with $d(u)=d(v)=5$ and
$d_G(u,v)\ge3$.  Thus $S$ is independent, no outside vertex is
adjacent to both members, and $H=G-S$ is smaller and nonempty.
Choose two nonadjacent neighbours $x,y$ of $u$; they exist since its
five neighbours cannot induce a planar $K_5$.  Properly six-colour the
planar $H-\{x,y\}$, and give both $x,y$ colour seven.  This is a
proper colouring of $H$ in which the neighbours of $u$ use at most
four colours, leaving at least three available.  The restriction
colourings differ at $w$, so Lemma~\ref{lem:independent-hc} gives
an $\alpha$--$\beta$ Hamilton path.

For $k\ge8$, planar five-degeneracy and
Theorem~\ref{thm:degenerate-hc} suffice.  For sharpness, colour the
six antipodal vertex pairs of the icosahedron with six different
colours.  Each vertex sees all other five colours, so this is a frozen
proper six-colouring; its recolouring graph has more than one state and
is disconnected.  Thus both thresholds are seven.
\end{proof}

Planar five-degeneracy and Theorem~\ref{thm:degenerate-vp} also give
vertex-pancyclicity for $k\ge8$.

\subsection{Triangle-free and bipartite planar graphs}

The critical five-colour step needs an adjacent-pair fibre, which is
not a product of complete graphs.  We give its finite input separately
from the arbitrary-size lifting argument.

For $A,B\subseteq\{0,1,2,3,4\}$, put
\[
 F(A,B)=\{(a,b)\in A\times B:a\ne b\},
\]
with adjacency given by changing one coordinate.  When $|A|,|B|$ are
three or four, exchange the coordinates if necessary so that
$|A|\le|B|$.  There are exactly seven types
$(|A|,|B|,|A\cap B|)$.

\begin{lemma}[The seven adjacent-pair types]
\label{lem:pair-five-endpoints}
The five types $331,342,343,443,444$ are Hamilton-connected.
The type $333$ is a $6$-cycle, so its feasible Hamilton-path endpoint
pairs are exactly its six edges.  For the normalized type $332$,
$A=\{0,1,2\}$ and $B=\{0,1,3\}$, all distinct endpoint pairs are
feasible except $\{03,21\}$ and $\{13,20\}$, where $ab$ denotes
$(a,b)$.
\end{lemma}

\begin{proof}
The explicit words below prove every positive assertion: each word
lists the complete state set once, with consecutive words differing in
one coordinate.  Apply all colour permutations preserving the listed
pair of lists, the coordinate exchange when it preserves the type, and
path reversal.  The resulting endpoint pairs have the counts in the
following table, covering all claimed positive pairs.
\[
\begin{array}{c|r|r|r}
 \text{type}&|V(F)|&\text{words}&\text{feasible unordered pairs}\\\hline
 331&8&6&28\\
 332&7&7&19\\
 333&6&1&6\\
 342&10&16&45\\
 343&9&8&36\\
 443&13&11&78\\
 444&12&4&66
\end{array}
\]
For $332$, the two excluded pairs are the neighbour pairs of the
degree-two states $01$ and $10$, respectively.  Choosing such a pair
as spanning-path ends would force the entire path to have only three
vertices.  For $333$, inspection gives the stated $6$-cycle.  These
supply exact negative witnesses, rather than search failures.
\end{proof}

\paragraph{The 53 endpoint words.}
For $331$, use $A=\{0,1,2\}$ and $B=\{0,3,4\}$:
\begin{verbatim}
03 13 23 24 20 10 14 04
03 23 20 24 04 14 13 10
03 23 20 24 04 14 10 13
03 23 13 10 20 24 04 14
13 10 20 24 23 03 04 14
13 14 10 20 23 03 04 24
\end{verbatim}
For $332$, use $A=\{0,1,2\}$ and $B=\{0,1,3\}$:
\begin{verbatim}
01 21 23 20 10 13 03
01 21 20 23 03 13 10
01 03 23 21 20 10 13
01 21 20 10 13 03 23
03 01 21 23 20 10 13
03 01 21 23 13 10 20
03 01 21 20 10 13 23
\end{verbatim}
For $333$, use $A=B=\{0,1,2\}$:
\begin{verbatim}
01 21 20 10 12 02
\end{verbatim}
For $342$, use $A=\{0,1,2\}$ and $B=\{0,1,3,4\}$:
\begin{verbatim}
01 21 24 20 23 13 10 14 04 03
01 21 23 20 24 14 04 03 13 10
01 21 24 23 20 10 14 04 03 13
01 21 23 24 14 04 03 13 10 20
01 04 14 24 23 03 13 10 20 21
01 21 24 20 10 13 14 04 03 23
03 23 24 14 13 10 20 21 01 04
03 23 24 20 21 01 04 14 10 13
03 13 10 20 24 23 21 01 04 14
03 23 24 21 01 04 14 13 10 20
03 13 23 24 20 10 14 04 01 21
03 01 21 20 24 04 14 10 13 23
03 23 13 14 10 20 21 01 04 24
20 24 23 13 10 14 04 03 01 21
20 10 13 14 04 24 21 01 03 23
23 21 20 10 14 13 03 01 04 24
\end{verbatim}
For $343$, use $A=\{0,1,2\}$ and $B=\{0,1,2,3\}$:
\begin{verbatim}
01 21 23 20 10 12 13 03 02
01 21 20 23 13 10 12 02 03
01 21 20 23 13 03 02 12 10
01 21 23 20 10 13 03 02 12
01 21 23 20 10 12 02 03 13
01 02 03 23 13 12 10 20 21
01 21 20 10 13 12 02 03 23
03 23 20 21 01 02 12 10 13
\end{verbatim}
For $443$, use $A=\{0,1,2,3\}$ and $B=\{0,1,2,4\}$:
\begin{verbatim}
01 31 34 24 21 20 30 32 12 10 14 04 02
01 21 31 34 32 30 20 24 14 10 12 02 04
01 21 31 34 32 30 20 24 14 04 02 12 10
01 31 32 30 34 24 21 20 10 14 04 02 12
01 31 32 30 34 24 21 20 10 12 02 04 14
01 21 20 30 32 31 34 14 10 12 02 04 24
01 21 31 32 30 20 24 14 10 12 02 04 34
04 24 34 32 31 30 20 21 01 02 12 10 14
04 34 14 24 20 21 01 31 32 02 12 10 30
04 34 32 30 20 21 24 14 10 12 02 01 31
04 24 20 30 32 31 21 01 02 12 10 14 34
\end{verbatim}
For $444$, use $A=B=\{0,1,2,3\}$:
\begin{verbatim}
01 31 21 23 20 30 32 12 10 13 03 02
01 21 31 32 30 20 23 13 03 02 12 10
01 21 31 32 30 10 20 23 13 03 02 12
01 21 31 32 30 20 10 13 12 02 03 23
\end{verbatim}

\begin{lemma}[Adjacent-pair interfaces and reset]
\label{lem:pair-five-reset}
Arrange full fibres $F(A_i,B_i)$ along a base path.  Each list has
size three or four in a fixed five-colour palette.  At each step either
both lists stay unchanged, or one stays unchanged and the other is
replaced by a three-element list differing in one colour, or a
three-element list and a containing four-element list are interchanged.
Join the copies of equal assignments at consecutive layers.  If an
internal layer has $|A_r|=|B_r|=4$, there is a spanning path between
any prescribed first-layer and last-layer states.
\end{lemma}

\begin{proof}
First consider one allowed interface between $F(A,B)$ and $F(A',B')$,
and write $J$ for their common assignment set.  It has at least four
states.  For every $p\in F(A,B)$, there is a $q\in J$ such that the
first fibre has a $p$--$q$ Hamilton path.  In a Hamilton-connected
type choose any $q\ne p$.  In type $332$, each $p$ forbids only
itself and at most one additional endpoint, so again a choice exists.
In type $333$, $J$ is either its whole $6$-cycle or that cycle minus
a row or column of two adjacent states.  A vertex's two cycle
neighbours are not adjacent to one another, so at least one remains in
$J$.  The Hamilton path obtained by deleting their connecting cycle
edge has the desired endpoints.  The same property holds at the
reverse interface.

Use this property to propagate nonempty sets of reachable exits from
the prescribed first state.  The internal $44$ layer is
Hamilton-connected, and its outgoing interface has at least nine states.
Any reachable entrance $p$ can therefore reach all interface states
except possibly $p$, yielding at least two distinct reachable exits.
The next layer still has one four-element list and is one of
$342,343,443,444$, hence is Hamilton-connected.  Two distinct reachable
entrances let every state of that layer be an exit.  If it is the last
layer, this includes the prescribed endpoint.  Otherwise its entire
outgoing interface is reachable.

Thereafter this property propagates: for a specified exit $q$ in the
current layer, apply the one-interface property in reverse to find a
Hamilton-path entrance in its fully reachable incoming interface.
The last prescribed endpoint is therefore reachable.  Backtracking
gives one Hamilton path in every full fibre, joined by actual cross
edges, with no repetition or omission.
\end{proof}

\begin{theorem}[Triangle-free planar Hamilton-connectedness]
\label{thm:trianglefree-hc}
If $G$ is a nonempty triangle-free planar graph, then $\Rk(G)$ is
Hamilton-connected for every $k\ge5$.  For the class of nonempty
triangle-free planar graphs,
$h=\hc=5$.
\end{theorem}

\begin{proof}
For $k\ge6$, use three-degeneracy and
Theorem~\ref{thm:degenerate-hc}.  For $k=5$, induct on the order of
$G$, permitting disconnected graphs.  Orders at most three are handled
by two-degeneracy and the same theorem.  At larger orders a vertex $v$
of degree at most two has a non-neighbour.  The induction hypothesis
gives a Hamilton-connected base of at least five states, thus of minimum
degree at least three; Lemma~\ref{lem:single-hc} applies.

Assume now $\delta(G)\ge3$.  Let $L$ be its degree-three vertices
and $s=|L|$.  The triangle-free planar edge bound gives
\[
 s\ge8+\sum_{d_G(z)\ge4}(d_G(z)-4).
\]
Given distinct $\alpha,\beta$, choose $w$ on which they differ.
There are two exhaustive cases.

If $L\setminus\{w\}$ contains adjacent $u,v$, denote their other
neighbours by $u_1,u_2$ and $v_1,v_2$.  The four actual outside
vertices are distinct: a common neighbour of $u$ and $v$ would form
a triangle.  In $H=G-\{u,v\}$ the restriction endpoints still differ
at $w$, so induction gives a base Hamilton path between them.  Its
fibre over $f$ is
\[
 F(A_f,B_f),\qquad
 A_f=[5]\setminus\{f(u_1),f(u_2)\},\quad
 B_f=[5]\setminus\{f(v_1),f(v_2)\}.
\]
The lists have size three or four.  One actual base move affects at
most one of them and gives precisely an interface allowed in
Lemma~\ref{lem:pair-five-reset}.  Edges among the four outside
vertices restrict the base colourings, but do not alter this statement.

There is a base colouring with both lists four-element: use
Gr\"otzsch's theorem~\cite{Grotzsch59} to three-colour
$H-\{u_1,u_2,v_1,v_2\}$, then give $u_1,u_2$ colour four and
$v_1,v_2$ colour five.  Each pair is independent by triangle-freeness,
and edges across the pairs have distinct colours.  At least five global
colour permutations give distinct double-four-list base states.  Since
the Hamilton path has only two endpoints, one is internal.  Apply
Lemma~\ref{lem:pair-five-reset} with the specified local endpoint
states to complete the restoration.

Otherwise $S=L\setminus\{w\}$ is independent and has at least
$s-1\ge7$ vertices.  For any outside vertex $z$ of degree at least
four, the displayed Euler budget implies $d_G(z)\le s-4<|S|$.
The only possible degree-three outside vertex is $w$, whose degree is
also less than $|S|$.  Thus no outside vertex sees all of $S$.
Choose $u\in S$ and two of its neighbours $x,y$.  They are
nonadjacent.  Three-colour $H-\{x,y\}$ by Gr\"otzsch's theorem
and give both $x,y$ colour five.  At most two colours then occur on the
three neighbours of $u$, so $u$ has at least three available colours.
All conditions of Lemma~\ref{lem:independent-hc} hold, and $w$ is
retained; it gives the desired endpoint path.

For sharpness, on the cube $Q_3$ view the colours as
$\mathbb Z_2^2$ and set
\[
 f(x,y,z)=x(1,0)+y(0,1)+z(1,1).
\]
The three neighbours of every vertex have all three other colours,
so the colouring is frozen and the complete four-colour recolouring
graph is disconnected.  The upper and lower bounds give both thresholds.
\end{proof}

\begin{corollary}
The class of nonempty bipartite planar graphs has $h=\hc=5$.
For every graph in this class, $\Rk(G)$ is vertex-pancyclic for
$k\ge6$.
\end{corollary}

\begin{proof}
Bipartite planar graphs are triangle-free, and the cube belongs to this
class.  The vertex-pancyclic assertion follows from three-degeneracy and
Theorem~\ref{thm:degenerate-vp}.
\end{proof}

The same vertex-pancyclic bound $k\ge6$ holds for all triangle-free
planar graphs.  This conclusion uses the degenerate theorem separately
from Hamilton-connectedness.

\subsection{Girth at least six and forests}

\begin{lemma}[Separated degree-two vertices with a marked vertex]
\label{lem:marked-two}
If $G$ is planar, $\girth(G)\ge6$, and $\delta(G)\ge2$, then for
any $w\in V(G)$ there are two degree-two vertices different from
$w$ whose distance is at least three.
\end{lemma}

\begin{proof}
Every component has a cycle, and the girth-six Euler bound gives at
least six degree-two vertices per component.  If there are at least two
components, choose such a vertex different from $w$ in each of two
components.  Suppose $G$ is connected and let $n_2$ be its number of
degree-two vertices.  From $2|E(G)|\le3|V(G)|-6$,
\[
 n_2\ge6+\sum_{d_G(z)\ge3}(d_G(z)-3).
\]
Let $L$ consist of all degree-two vertices except $w$; then
$|L|\ge n_2-1\ge5$.  Suppose all pairs in $L$ have distance at
most two.

If $u,v\in L$ are adjacent, let $a,b$ be their respective other
neighbours.  The absence of cycles of length at most five implies that
the intersection of the two closed distance-two balls is exactly
$\{u,v,a,b\}$, and that $d_G(a,b)=3$.  Thus at most three of
these vertices can lie in $L$, a contradiction.

Otherwise $L$ is independent.  Its two-element neighbourhoods are
pairwise intersecting; they are distinct because there is no four-cycle.
At least four distinct pairwise intersecting two-element sets have a
common element: if $\{a,b\},\{a,c\}$ occur and a third set omits
$a$, it must be $\{b,c\}$, after which no fourth distinct set is
possible.  Hence some vertex $z$ is adjacent to all of $L$.
Then $d_G(z)\ge n_2-1\ge5$, whereas the Euler inequality gives
\[
 n_2\ge6+d_G(z)-3\ge n_2+2,
\]
a contradiction.
\end{proof}

\begin{theorem}[Girth at least six]
\label{thm:girth6-hc}
For every finite integer $g\ge6$, the class of nonempty planar graphs
of girth at least $g$ has $h=\hc=4$.  Every graph $G$ in this
class has Hamilton-connected $\Rk(G)$ for $k\ge4$, and
vertex-pancyclic $\Rk(G)$ for $k\ge5$.
\end{theorem}

\begin{proof}
For $k\ge5$, use two-degeneracy and
Theorems~\ref{thm:degenerate-hc} and~\ref{thm:degenerate-vp}.
At $k=4$, induct on the order, starting with orders at most two,
which are forests and covered by the one-degenerate theorem.
If $|V(G)|\ge3$ and $d_G(v)\le1$, then $v$ has a non-neighbour.
The smaller nonempty base is Hamilton-connected by induction and has
at least four states, hence minimum degree at least three.  Restore
$v$ using Lemma~\ref{lem:single-hc}.

Otherwise apply Lemma~\ref{lem:marked-two} to a vertex $w$ at which
the prescribed endpoint colourings differ.  It gives separated
degree-two vertices $u,v$, both different from $w$.  Set
$S=\{u,v\}$ and $H=G-S$.  The set $S$ is independent, no outside
vertex sees both its members, and $H$ is nonempty.  The neighbours
$x,y$ of $u$ are nonadjacent.  The two-degenerate $H-\{x,y\}$
can be three-coloured; giving $x,y$ colour four leaves three colours
available at $u$.  Lemma~\ref{lem:independent-hc} now restores a
Hamilton path between the actual prescribed states.

For the lower bound, take a cycle of length at least $g$ divisible by
three and repeat the colours $0,1,2$ around it.  This is a frozen
proper three-colouring, so its complete three-colour recolouring graph
is disconnected.  Both class thresholds are therefore four.
\end{proof}

\begin{theorem}[Forests]
\label{thm:forest-hc}
The class of nonempty forests has $h=\hc=4$.  Every such forest
$F$ has Hamilton-connected and vertex-pancyclic $\Rk(F)$ for
$k\ge4$.
\end{theorem}

\begin{proof}
The upper bounds follow from one-degeneracy and
Theorems~\ref{thm:degenerate-hc} and~\ref{thm:degenerate-vp}.
For the Hamiltonian lower bound, $\R_3(P_3)$ has twelve
states.  Its six states using three distinct colours have degree two,
so every edge incident with them is forced in a spanning cycle.
For each of the three middle-vertex colours $b$, these forced edges
form a $4$-cycle
\[
 (a,b,a),(c,b,a),(c,b,c),(a,b,c),
 \qquad \{a,b,c\}=\{0,1,2\}.
\]
The three disjoint forced cycles prevent a Hamilton cycle.  For the
Hamilton-connected lower bound, the subgraph $\R_3(K_2)=C_6$
is already not Hamilton-connected.
\end{proof}

\subsection{The single graph \texorpdfstring{$C_5$}{C5}}

\begin{theorem}[The five-cycle]
\label{thm:c5}
The individual graph $C_5$ has $h(C_5)=\hc(C_5)=4$.
Its four-colour recolouring graph is Hamilton-connected, its
three-colour recolouring graph is the disjoint union of two $15$-cycles,
and $\Rk(C_5)$ is Hamilton-connected and vertex-pancyclic for
$k\ge5$.
\end{theorem}

\begin{proof}
The four-colour assertion is a complete finite certificate.  The state
set consists of the 240 proper words of length five over four colour
labels.  Under colour permutation and the dihedral symmetries of the
cycle, every state is equivalent to one of $01012$ and $01023$.
For each of these two initial states the certificate records a Hamilton
path to each of the other 239 states.  Its 478 paths each contain all
240 states exactly once, have the required endpoints, and have a
single-coordinate change at each of the 114242 steps.  The group
actions preserve properness and adjacency, so these paths cover every
ordered pair of distinct endpoint states.  This establishes precisely
the complete finite graph $\R_4(C_5)$.

Direct enumeration at three colours gives 30 proper states, all of
degree two, with two components of order fifteen.  Hence that graph is
two $15$-cycles and is not Hamiltonian or Hamilton-connected.
Finally, $C_5$ is two-degenerate, so the separate assertions for
$k\ge5$ follow from Theorems~\ref{thm:degenerate-hc}
and~\ref{thm:degenerate-vp}.  The finite certificate makes no claim
about other girth-five graphs.
\end{proof}


\section{Hamiltonian four-colour subclasses at girth five}
\label{sec:four-colour-subclasses}

The results in this section are Hamiltonian results.  They do not assert
Hamilton-connectedness at four colours.  The local finite inputs are recorded below.  Their finite parts classify
labelled lists; the inductions are valid for an arbitrary number of
vertices and are not extrapolations from finite state spaces.







\begin{lemma}[Subdivided-star lift]
\label{lem:subdivided-star-four-subclass}
Let $S=\{u_1,\ldots,u_s\}$ be independent, $s\ge3$, and suppose
$N(u_i)=\{w,x_i\}$, where $N(w)=S$, $w\notin S$, and the $x_i$ are distinct and belong to
$J=G-(S\cup\{w\})$.  If $\R_4(J)$ is Hamiltonian and some proper colouring
of $J$ gives the same colour to two distinct $x_i$, then $\R_4(G)$ is
Hamiltonian.
\end{lemma}

\begin{proof}
Take a Hamilton cycle $C$ of $\R_4(J)$.  Choose three distinct edges of
$C$.  At each edge at most one $x_i$ changes, so an unchanged $x_j$ supplies
a colour $a_r$ at the $r$th chosen edge.  Choose a tree on the four colours
whose three edges contain $a_1,a_2,a_3$, respectively (a star works when
there is one repeated label; if the labels are all distinct, use the star
at the fourth colour).  Make four copies of $C$, indexed by the colour of
$w$, and splice the corresponding copies along the three selected edges
according to that colour tree.  The result is a Hamilton cycle $B$ of
$\R_4(G-S)$ with the following control: whenever $w$ changes from $a$ to $b$
along $B$, some $x_j$ has colour $a$ or $b$.

Over a base colouring $f$ the $u_i$-lists are
\[
 A_i(f)=[4]\setminus\{f(w),f(x_i)\}.
\]
They have size at least two.  At a move in $J$, one list changes and another
coordinate is unchanged.  At a move of $w$, all intersections are
$[4]\setminus\{a,b,f(x_i)\}$, and the control property leaves a coordinate
with a two-element intersection.  Thus the product-fibre splice applies at
every step.  At the base colouring with $f(x_p)=f(x_q)$, give $w$ that same
colour; the two corresponding lists have size three, so both interfaces at
this fibre contain two distinct splice edges.  The resulting cycle traverses
all extensions and proves the lemma.
\end{proof}

\begin{lemma}[Five-face identity]
\label{lem:five-face-identity-subclass}
Let $C$ be a connected plane graph of girth at least five and minimum degree
at least two.  Let $p$ be the number of $5$-faces,
$Q=\sum_{|F|\ge7}(|F|-6)$, $s=|\{v:d(v)=2\}|$, and
$D=\sum_v(d(v)-3)$ over vertices of degree at least three.  Then
\[
 p=12+Q+2D-2s.
\]
Every $5$-face is a simple $5$-cycle, and distinct faces give distinct
cycles unless $C=C_5$.
\end{lemma}

\begin{proof}
A facial walk has no immediate reversal and contains a cycle, hence has
length at least five.  Counting face lengths gives
$2e=6f-p+Q$, while the degree sum is $2e=3n-s+D$; with $n-e+f=2$,
subtraction yields the identity.  A repeated vertex on a facial walk of
length five would give a shorter cycle.  If one $5$-cycle bounds two faces,
all its vertices have degree two and connectedness gives $C=C_5$.
\end{proof}

\begin{corollary}[Two-degeneracy from seven pentagons]
\label{cor:seven-pentagons-two-degenerate}
A planar graph of girth at least five with at most seven distinct $5$-cycles
is two-degenerate.
\end{corollary}

\begin{proof}
Otherwise a connected subgraph has minimum degree at least three.  In its
plane embedding the identity has $s=0$, and therefore at least twelve
$5$-faces, which are distinct $5$-cycles.  This contradicts the hypothesis.
\end{proof}

\begin{theorem}[At most seven distinct five-cycles]
\label{thm:seven-pentagons}
Let $G$ be a nonempty planar graph with girth at least five and at most
seven distinct $5$-cycles.  Then $\Rk(G)$ is Hamiltonian for every $k\ge4$;
therefore the Hamiltonian threshold of this class is four.  No
four-colour Hamilton-connected conclusion is claimed.
\end{theorem}

\begin{proof}
For $k=4$ use induction on $|V(G)|$, allowing disconnected graphs.  The
one-vertex graph gives $K_4$.  A vertex of degree at most one is restored by
Lemma~\ref{lem:single-lift}.  If two degree-two vertices are adjacent, use
Lemma~\ref{lem:adjacent-two-four}.  Otherwise the degree-two
vertices form an independent set $S$.  In a component containing a cycle,
Lemma~\ref{lem:five-face-identity-subclass} and the bound $p\le7$ give
$|S|\ge3$; a tree component has a leaf and was already reduced.

If no vertex is adjacent to all of $S$, every vertex outside $S$ has at most
$|S|-1$ neighbours in $S$.  Choose $u\in S$ with neighbours $x,y$.  They are
nonadjacent by the girth condition.  Three-colour $G-S-\{x,y\}$ greedily
using Corollary~\ref{cor:seven-pentagons-two-degenerate}, and colour both
$x,y$ with a fourth colour.  This gives a base colouring with three colours
available at $u$.  Lemma~\ref{lem:independent-restoration} restores $S$.

Suppose that a vertex $w$ is adjacent to all of $S$.  In the component under
consideration write $d(w)=|S|+t$.  The identity and
$D\ge d(w)-3$ give $p\ge6+Q+2t$, so $p\le7$ forces $t=0$ and
$N(w)=S$.  Write $N(u_i)=\{w,x_i\}$.  The $x_i$ are distinct, and two are
nonadjacent since $|S|\ge3$ and the graph has no triangle.  Put
$J=G-(S\cup\{w\})$.  Greedy three-colouring of $J-\{x_p,x_q\}$ followed by
assigning the fourth colour to $x_p,x_q$ supplies the hypothesis of
Lemma~\ref{lem:subdivided-star-four-subclass}; induction on $J$ completes
this case.  For $k\ge5$, apply the degenerate Hamiltonian theorem to the
corollary above (and use the one-vertex lift at each deletion); this is a
separate argument and assumes no monotonicity in $k$.
\end{proof}

\begin{lemma}[Adjacent degree-two restoration]
\label{lem:adjacent-two-four}
Let $u,v$ be adjacent degree-two vertices, with
$N(u)=\{v,x\}$ and $N(v)=\{u,y\}$, where $x\ne y$.  If
$\R_4(G-\{u,v\})$ is Hamiltonian, then $\R_4(G)$ is Hamiltonian.
\end{lemma}
\begin{proof}
For a base colouring $f$, the extension fibre is
$F(A,B)$ with $A=[4]\setminus\{f(x)\}$ and
$B=[4]\setminus\{f(y)\}$.  If $A=B=\{a,b,c\}$, use
\[
(ab,ac,bc,ba,ca,cb).
\]
If $A=\{a,b,c\}$ and $B=\{a,b,d\}$, use either
\[
(ab,ad,bd,ba,ca,cd,cb),\qquad
(ab,ad,cd,bd,ba,ca,cb).
\]
A one-colour list change has two common labelled internal edges between
any old and new selected cycles; in the equal case use the two unchanged
rows, and in the unequal case inspect the two displayed words.  Cut an
edge of a base Hamilton cycle and splice successive fibres along these
two common edges.  Every splice joins two disjoint cycles and the fibre
states are partitioned, so the resulting cycle is spanning.
\end{proof}

\subsection{Two degree-two reductions}
\begin{lemma}[Separated degree-two vertices]\label{lem:separated-two-g5}
Let $G$ be planar of girth at least five and let $u,v$ be nonadjacent
vertices of degree two with $N(u)\cap N(v)=\varnothing$.  If
$H=G-\{u,v\}$ has Hamiltonian $\R_4(H)$, then $\R_4(G)$ is Hamiltonian.
\end{lemma}
\begin{proof}
Set $S=\{u,v\}$.  Every vertex of $H$ has at most one neighbour in $S$.
Let $x,y$ be the neighbours of $u$; they are distinct and nonadjacent.
The planar triangle-free graph $H-\{x,y\}$ admits a proper three-colouring
by Gr\"otzsch's theorem~\cite{Grotzsch59} (the empty graph has its empty
colouring).  Use colours $0,1,2$ there and give both $x,y$ colour $3$.
This is a proper colouring of $H$ leaving three colours available at $u$.
The hypotheses of Lemma~\ref{lem:independent-restoration} hold with
$k=4$ and $|S|=2$, so that lemma supplies the required Hamilton cycle.
\end{proof}

\subsection{\texorpdfstring{Restoring a $(2,3,2)$ path}{Restoring a (2,3,2) path}}
The following reduction uses only a Hamilton cycle of the deleted graph.
It processes an entire constant-boundary run together with its two
neighbouring nonconstant fibres, retaining the edges needed for the next
splice.  The three external neighbours remain distinct actual vertices.
\begin{proposition}\label{prop:badrun}
If $F=\R_3(P_3)$, then $F\mathbin{\square}P_\ell$ is Hamiltonian
for every $\ell\ge2$.
\end{proposition}
\begin{proof}
Partition $F$ into its three centre-colour squares $Q_b$.
Order each square as $(q_0,q_1,q_2,q_3)$, where $q_0,q_2$ are
its monochromatic-endpoint vertices. On $Q_b\square P_2$ use
\[
 (0,q_0),(1,q_0),(1,q_1),(0,q_1),
 (0,q_2),(1,q_2),(1,q_3),(0,q_3).
\]
This cube cycle contains all four edges between layers 0 and 1.
For larger $\ell$, append each subsequent square by an ordinary
edge splice using an internal edge of the newest layer. The initial
newest layer has two internal edges; each appended layer retains
three. This leaves all four original interlayer edges intact.

The three centre-colour squares of $F$ are connected by three edges
between monochromatic-endpoint vertices, one for each pair of
centre colours. Choose a tree of two of these edges. Each chosen
edge $xy$ produces a square with vertices $(0,x),(1,x),(1,y),(0,y)$.
Delete its two interlayer edges and add its two copies of $xy$,
joining the corresponding strip cycles. The two selected interlayer
edges at the middle strip are distinct. Two such splices give a
single spanning cycle, for arbitrary $\ell\ge2$.
\end{proof}

Let $x,y,z$ be a path of degrees $2,3,2$.  Its external neighbours, in
that order, are $p,q,r$.  Assume they are distinct and outside the path.
For a bottom colouring with boundary $t=(f(p),f(q),f(r))$, write
\[
 B_t=\{(u,v,w):u\ne v,\ v\ne w,\ u\ne t_0,\ v\ne t_1,
                  \ w\ne t_2\},\qquad u,v,w\in\{0,1,2,3\}.
\]
A boundary is \emph{bad} if all three entries agree, and \emph{good}
otherwise. Up to colour permutations and path reversal the good types
are $001,010,012$. Their orders are $6+4+4=14$, $9+4+4=17$, and
$6+6+4=16$, respectively. The bad type $000$ has order $12$.

\subsection{Good cycles with safe connecting edges}
In the following cyclic words the coordinates are $uvw$:
\[
\begin{array}{c|l}
001&120\ 123\ 323\ 313\ 213\ 210\ 230\ 130\ 132\ 232\ 212\ 312\ 310\ 320\\[2pt]
010&101\ 102\ 103\ 203\ 303\ 323\ 123\ 121\ 321\\
   &301\ 302\ 202\ 201\ 231\ 232\ 132\ 131\\[2pt]
012&101\ 121\ 120\ 320\ 321\ 323\ 123\ 103\\
   &203\ 303\ 301\ 201\ 231\ 230\ 130\ 131
\end{array}                                                     \tag{8}
\]
Call these cycles $D_E,D_M,D_A$, and allow all their colour images and
path reversals. Legality, cyclic adjacency, distinctness, and the counts
above verify that each is a Hamilton cycle of its fibre.

For a bad boundary $aaa$, partition $B_{aaa}$ into squares $Q_c$ indexed
by $c\ne a$, with fixed middle coordinate $c$. A good neighbour of
$aaa$ is obtained by changing position $j$ from $a$ to $b\ne a$.
If $j$ is an end position, then for each $c\notin\{a,b\}$ there is
exactly one edge common to the good fibre and $Q_c$. Let $P_t$ consist
of these two edges. For middle-change or all-distinct boundaries set
$P_t=\varnothing$. A good-cycle edge outside $P_t$ is called \emph{safe}.
Directly from (8), every allowed cycle has the following properties:
\begin{enumerate}
\item For an end change, it contains the unique edge in each
      $Q_c$, $c\notin\{a,b\}$, and two opposite edges of $Q_b$.
\item For a middle change it has no vertex in $Q_b$ and contains
      three edges of each of the other two squares.
\end{enumerate}
For $t=001$, the two special edges are $123$--$323$ and $132$--$232$;
the two edges in $Q_1$ are $313$--$213$ and $212$--$312$.
For $t=010$, the retained three-edge paths are
$323,123,121,321$ and $231,232,132,131$.
These descriptions also prove the assertions for every symmetry image.

\begin{lemma}[Safe common edges]\label{lem:pathsafe}
For adjacent good boundaries $t,s$, every allowed cycle at $t$ and
every allowed cycle at $s$ have at least two common edges outside
$P_t\cup P_s$.
\end{lemma}
\begin{proof}
Number source edges in (8) cyclically from one. The following explicit
witness table supplies two safe common source-edge indices in every
case. Target columns are defined as follows.

For an end-change target, write $a$ for its repeated colour, $b$ for
its singleton, and $h_0<h_1$ for the unused colours. There is one
allowed edge set: apply $p=(a,b,h_0,h_1)$ to $D_E$, reversing the path
if the singleton is at the left end. Swapping the two unused colours
preserves this edge set.
For a middle-change target $(a,b,a)$ there are four images of $D_M$:
column $\epsilon\eta$ uses $p=(a,b,h_\eta,h_{1-\eta})$ and reverses
the path exactly when $\epsilon=1$. For an all-distinct target
$(a,b,c)$ with unused colour $d$, columns 1 and 2 are respectively
$p=(a,b,c,d)$ applied to $D_A$, and $p=(c,b,a,d)$ followed by path
reversal. These exhaust the images in each case.
\begin{center}\small
\begin{tabular}{cc|cccc}
$t$ & $s$ & $1$ / $00$ & $2$ / $01$ & $10$ & $11$\\\hline
$001$ & $002$ & $1,3$ & -- & -- & --\\
$001$ & $003$ & $6,8$ & -- & -- & --\\
$001$ & $011$ & $1,7$ & -- & -- & --\\
$001$ & $021$ & $4,7$ & $4,5$ & -- & --\\
$001$ & $031$ & $1,5$ & $1,3$ & -- & --\\
$001$ & $101$ & $3,4$ & $4,5$ & $3,12$ & $5,10$\\
$001$ & $201$ & $8,12$ & $1,8$ & -- & --\\
$001$ & $301$ & $1,5$ & $1,5$ & -- & --\\
$010$ & $012$ & $3,4$ & $3,6$ & -- & --\\
$010$ & $013$ & $9,11$ & $8,10$ & -- & --\\
$010$ & $020$ & $1,3$ & $2,14$ & $3,10$ & $2,4$\\
$010$ & $030$ & $5,7$ & $2,6$ & $1,3$ & $1,2$\\
$010$ & $210$ & $1,10$ & $2,6$ & -- & --\\
$010$ & $310$ & $3,7$ & $2,13$ & -- & --\\
$012$ & $013$ & $1,2$ & $3,4$ & -- & --\\
$012$ & $032$ & $3,8$ & $2,4$ & -- & --\\
$012$ & $312$ & $8,12$ & $1,7$ & -- & --\\
\end{tabular}
\end{center}

To prove completeness, order the types as $E<M<A$ and reverse the
boundary transition so that its source has no larger type than its
target. Normalize its chosen source cycle to (8) by a colour permutation
and possibly reversal. There are eight eligible targets from $001$,
six from $010$, and three from $012$, exactly the seventeen rows above.
The column rules exhaust every target cycle. Each of the thirty-seven
entries is verified by substituting its stated colour map and reversal
in (8), and excluding the two explicitly described special edges where
applicable. Thus the table is a complete finite analytic certificate.
The transformations preserve special edges, completing the proof.
\end{proof}

\subsection{A bad block with prescribed safe edges left intact}
\begin{lemma}[Bad-block lifting]\label{lem:badblock}
Consider consecutive fibres with boundaries
\[
             t,\underbrace{aaa,\ldots,aaa}_{\ell\text{ times}},s,
             \qquad \ell\ge1,
\]
where $t,s$ are good and each differs from $aaa$ in one coordinate.
When $\ell=1$ suppose $t\ne s$. Choose any allowed good cycles at
the two ends, and reserve at most one safe edge in each. The layered
graph on these fibres has a Hamilton cycle retaining both reserved
edges. Only consecutive-layer matching edges are used.
\end{lemma}
\begin{proof}
Call a common edge with $Q_c$ \emph{available} if it belongs to the
chosen good cycle and is not its reserved edge. An end-change good
cycle has available edges in all three squares: the two singleton
edges are special, hence cannot have been reserved, and at least one
of its two $Q_b$ edges remains. A middle-change good cycle has at least
two available edges in each of its two supported squares.

First suppose $\ell=1$. Start with five disjoint cycles: the two good
cycles and the three bad squares. Choose a bridge square with two
different available edges, one for each good cycle; attach the other
two squares as leaves to good cycles. The following cases guarantee
such a choice. Denote the changed colours by $b_L,b_R$.
\begin{enumerate}
\item Both changes are at ends. Choose $c\notin\{a,b_L,b_R\}$.
      Both singleton edges in $Q_c$ are available. They are different:
      if the changed end is the same, $t\ne s$ implies $b_L\ne b_R$,
      and the edges are opposite in $Q_c$; if the changed ends differ,
      the edges change different endpoint coordinates. Either good
      cycle can receive either remaining square.
\item Both changes are at the middle. Here $b_L\ne b_R$. Use the
      remaining colour $c$ for the bridge square. Each side has at
      least two available edges there, so different edges can be
      selected. Attach $Q_{b_L}$ to the right and $Q_{b_R}$ to the left.
\item One change is at the middle, with colour $b_M$, and the other
      at an end, with colour $b_E$. Choose
      $c\notin\{a,b_M,b_E\}$. The end's singleton edge in $Q_c$ is
      available, and the middle side has an available edge different
      from it. Attach $Q_{b_M}$ to the end-change cycle, and the last
      square to the middle-change cycle. This also covers $b_M=b_E$:
      the end-change cycle retains at least one of its two $Q_{b_M}$
      edges.
\end{enumerate}
The four selected incidences form a tree on the five initial cycles.
Cuts on a good cycle occur at different middle colours, hence are
distinct; the two bridge-square cuts were chosen distinct. Along each
tree edge remove the corresponding internal edge from each cycle and
add the two matching edges. Successively adding a new tree component
merges two cycles into one. All reserved edges survive.

Now let $\ell\ge2$. Each good end supports at least two squares, so
choose an available left edge in $Q_c$ and an available right edge in
$Q_d$ with $c\ne d$. Build the three square-strip cycles as in
Proposition~\ref{prop:badrun}. The cube formula there permits either
perfect matching of a square to be used in the first layer, by rotating
its cyclic order; the last layer of a two-layer strip uses the other
matching. Consequently its first layer can retain any specified square
edge. For a two-layer strip the same is true at its last layer; our two
requirements are on different strips and are independent.

If $\ell\ge3$, a last-layer requirement is met by appending each new
square with a splice edge different from the required edge. There are
initially two choices and subsequently three in the current last
layer, so this is always possible. The new last square retains all
edges except the splice edge. First-layer horizontal edges and all
four original vertical edges survive these extensions.

Join the three strips using the two vertical-edge switches in
Proposition~\ref{prop:badrun}; these do not affect either required
horizontal edge. Finally join the left and right good cycles to their
selected strip edges by two ordinary edge splices. The result is one
Hamilton cycle. The reserved edges were never cut. This proves the
claim for every length, rather than just a bounded set of runs.
\end{proof}

\subsection{The complete path-lifting induction}
\begin{lemma}[$(2,3,2)$-path lifting]\label{lem:path232}
Suppose the three external neighbours of a degree-$(2,3,2)$ path are
distinct and outside it. If the reconfiguration graph of the deleted
graph is Hamiltonian at four colours, then so is $\R_4(G)$.
\end{lemma}
\begin{proof}
Fix an arbitrary bottom Hamilton cycle. There must be a good vertex
on it. Otherwise every boundary is constant; consecutive constant
triples can only be equal, as at most one coordinate changes. The
external colour would therefore be fixed throughout the cycle. A
global colour permutation of a proper colouring supplies another
proper colouring with a different external colour, a contradiction
to Hamilton coverage. Thus even the startup uses only Hamiltonicity.

Every good bottom vertex has at most one bad neighbour in the entire
bottom reconfiguration graph: a two-colour boundary has exactly one
possible coordinate change making it constant, and an all-distinct
boundary has none. The full neighbouring colouring, if proper, is
uniquely determined by that change. Hence no good vertex on the
Hamilton cycle can lie between two bad vertices. In particular there
is an edge joining two good vertices.

Cut such an edge. In the resulting bottom Hamilton path, group each
maximal bad run together with its two neighbouring good vertices;
make each remaining good vertex a singleton block. These blocks are
disjoint and partition the path. A bad run has constant boundary,
and its two neighbours are distinct. If the run has length one their
boundaries are different: otherwise both neighbours are obtained from
the intervening colouring by the same one-vertex recolouring and
would be the same bottom vertex.

For each labelled good boundary fix one allowed cycle once and for
all. Consecutive blocks meet at a good--good boundary transition.
If the boundary changes, Lemma~\ref{lem:pathsafe} provides two common
safe edges. If it stays unchanged, the fixed cycles coincide and
have at least $14-2$ safe edges. Select one connecting edge for each
interblock interface in path order. If a singleton block already has
an incoming selected edge, choose a different outgoing one, which is
possible from the two candidates. In a nonsingleton block the two
selected edges are in its distinct end fibres, so cannot conflict.

Use Lemma~\ref{lem:badblock} to construct each nonsingleton block
cycle preserving its selected end edges; singleton blocks already
have their fixed cycles. Splice the resulting disjoint block cycles
along the selected interfaces. Each step joins two cycles into one,
and every later selected edge is still present. The one-block case
needs no final splice. All transitions are legal: a local step changes
one path vertex, and a matching cross-edge performs its actual bottom
recolouring while holding the local colouring fixed. Finally the fibres
are disjoint and exhaust all proper colourings. Thus the resulting
single cycle is Hamiltonian, including its closing edge.
\end{proof}



\begin{theorem}\label{thm:subcubic}
Let $G$ be a nonempty finite simple planar graph of girth at least five
and maximum degree at most three. Then $\Rk(G)$ is Hamiltonian for
every $k\ge4$. Disconnected graphs are allowed. The uniform Hamiltonian threshold
of this class is exactly four; four-colour Hamilton-connectedness is not asserted.
\end{theorem}

We prove a local restoration lemma using only Hamiltonicity of the
smaller recolouring graph, and then supply the planar structure.
The new literal words, pools and all 313 interface records appear in
Appendix~\ref{app:protected-pentagon-words}. They are explicit finite
identities in the proof, not references to solver verdicts.
The unrestricted planar girth-five four-colour question remains open.

\subsection{Listed pentagon factors}
\begin{lemma}\label{lem:pentagon-list-factor}
If every list on an induced pentagon is a subset of $Q$ of size at
least three, its list recolouring graph has such a factor.
\end{lemma}
\begin{proof}
When all lists have size three, record their omitted colours as $t$.
Up to the dihedral action and a common palette permutation, the
boundary types and factors are as follows.
\[
\begin{array}{c|r|l@{\qquad}c|r|l}
t&|F_t|&\text{factor}&t&|F_t|&\text{factor}\\\hline
00000&30&2C_{15}&00001&44&22K_2\\
00011&42&21K_2&00012&47&22K_2+K_3\\
00101&64&32K_2&00102&63&30K_2+K_3\\
00112&51&24K_2+K_3&00121&58&29K_2\\
00123&57&27K_2+K_3&01012&68&34K_2\\
01023&66&33K_2&&&
\end{array}
\]
The multiplicities $5,4+1,3+2,3+1+1,2+2+1,2+1+1+1$ yield,
respectively, $1,1,2,2,3,2$ position types on a pentagon.
Appendix~\ref{app:pentagon-factor-words} prints every factor: semicolons
separate blocks, and an unseparated line is a single block.
Successive words differ at one coordinate (cyclically for
a triangle or 15-cycle), and all words are legal and distinct.
For completeness the orders can be checked from
\[
 |F_t|=\operatorname{tr}(D_{t_0}A D_{t_1}A D_{t_2}A D_{t_3}A D_{t_4}A),
 \qquad A=J_4-I_4,
\]
where $D_c$ has diagonal zero at $c$ and one elsewhere. Expanding the
trace counts exactly the legal words. The 590 printed legal, distinct
words have these orders, proving coverage by explicit finite identities.
If a list is all of $Q$, fix the other four colours. The remaining
vertex has either two or three available colours, a clique fibre.
These fibres directly partition the state space.
\end{proof}

\subsection{The finite protected-edge identities}
Put $Q=\{0,1,2,3\}$. A boundary $t$ has five entries in $Q$, with at most
one entry $\ast$. Its local graph $F_t$ consists of the proper $C_5$
colourings $x$ satisfying $x_i\ne t_i$; $\ast$ excludes no colour.
In the certificates $\ast$ is encoded by 4.
The eleven ordinary types have 590 states in total.
The constant type has two components, each a 15-cycle; all other
ordinary types are connected. Eight have the printed Hamilton cycles,
and the two remaining types are $00001$ and $00112$.
With a full list, the eleven types, after rotation, reflection and
colour permutation, are
\[
0000\ast,\ 0001\ast,\ 0010\ast,\ 0011\ast,\ 0012\ast,\
0101\ast,\ 0102\ast,\ 0110\ast,\ 0112\ast,\ 0120\ast,\ 0123\ast.
\]
Their printed Hamilton cycles have respectively
$54,60,78,62,71,88,85,74,74,83,82$ states, a total of 811.
Thus there are 23 component types, including the two constant components.

For a Hamilton component $A$ with specified cycle $C_A$, designate an
ordered pool of eight edges of $C_A$, and let $K_A$ be the other cycle
edges. The first two pool edges will join consecutive copies of $C_A$;
the last six will join factor blocks.
For a hard component, specify cycles of $A\square P_2$ and
$A\square P_3$, and a set $D_A$ of internal edges occurring in
\emph{every} row of both seeds. We use new seeds for $00001$ with
$|D_A|=19$ and the earlier seeds for $00112$ with $|D_A|=18$.
Designate a seven-edge pool in $D_A$, and put $K_A=D_A\setminus
\text{pool}$. Its first edge joins successive seed blocks and its last
six join factor blocks. All pool edges are distinct.

Here is the exact normalization convention. Enumerate $s=0,\ldots,4$,
then $\epsilon=+1,-1$, then the permutations $p\in S_4$ in lexicographic
order. Take the first standard type
$n_i=p(t_{s+\epsilon i})$, retaining $\ast$. Restoration puts the
standard entry $x_i$ at position $s+\epsilon i$ with colour $p^{-1}(x_i)$.
In particular a standard boundary first normalizes by the identity.

\begin{lemma}[Protected-edge identities]\label{lem:pentagon-protection}
Let the source be a standard component $A$ of $F_t$, and change one
non-star boundary entry. Normalize the target by the preceding rule.
Every target component $B$ whose restored state set meets $A$ has an
edge in common with $K_A$ after restoration of $K_B$.
For an unchanged boundary use the same component and the same map.
\end{lemma}
\begin{proof}
The ordinary and starred orbit representatives cover
$4^5+5\cdot4^4=2304$ labelled boundaries. The displayed cyclic words
give their specified component cycles and the two hard seeds: their
entries are precisely the legal words, without repetition, and
successive entries, including last and first, differ in one coordinate.
Every listed common edge occurs in all five seed rows.
Removing the listed pools therefore defines the asserted protected sets.
For each standard source component, each non-star position, each
different new excluded colour and each overlapping normalized target
component, the appendix gives an actual edge in both protected sets.
There are 313 records. Checking membership of the two five-digit
endpoints in the specified cycles and outside their pools proves
each record directly. Applying the inverse coordinate and colour maps
gives every labelled case. This is a complete finite case proof;
no finite-length gluing experiment is used to infer a general rule.
\end{proof}

\subsection{Products retaining protection in every row}
\begin{lemma}\label{lem:listed-pentagon-degree}
For arbitrary lists contained in $Q$, a listed $C_5$ recolouring graph
has maximum degree at most seven.
\end{lemma}
\begin{proof}
In a proper four-colouring, a vertex has one possible move, or two when
its two neighbours have the same colour. Since $\alpha(C_5)=2$, each
colour occurs at most twice and at most two colours are repeated.
A repeated pair is the neighbour pair of exactly one vertex.
Hence there are at most two equal-neighbour corners and at most
$5+2=7$ moves. Restricting the lists only deletes moves.
\end{proof}

\begin{lemma}\label{lem:protected-product}
For any two of the 23 component types $A,B$, the graph $A\square B$
has a Hamilton cycle retaining a copy of every edge of $K_A$
in every physical $B$ row.
\end{lemma}
\begin{proof}
First construct $A\square P_r$ for every $r\ge2$.
When $C_A$ exists, join consecutive copies of it by alternating the
first two pool edges. At an intermediate row the two used edges are
distinct. Every join cuts one edge in two disjoint circles and joins
the resulting paths, giving a single circle.
When $A$ is hard, write $r$ as a sum of twos, preceded by one three if
$r$ is odd. Use the specified seeds on these intervals and join
successive intervals at their end rows using the first pool edge.
Every interval has at least two rows, so its two interfaces use
different physical rows. Both constructions retain $K_A$ and the six
remaining pool edges in every row.

Partition $B$ into ordered paths with at least two vertices: use its
edge/triangle factor in an ordinary nonconstant type, one 15-cycle
cut into a path in a constant component, or its printed Hamilton
cycle cut into a path in a starred type.
Contract these paths and choose a spanning tree in the quotient,
with an actual edge of $B$ witnessing each tree edge.
Every physical vertex has at least one neighbour within its factor
path. By Lemma~\ref{lem:listed-pentagon-degree}, it witnesses at most
six tree edges. Start with the path-product circles and join them
in parent-before-child order. The parent row has consumed at most
five of the six remaining pool edges, while the fresh child row has
all six. Choose a common unused edge for the two-edge switch.
Every step gives one circle on the blocks processed so far, and
never removes an edge of $K_A$. The final circle covers $A\square B$.
\end{proof}

\begin{lemma}\label{lem:pentagon-overlap}
For a fixed star position, or no star, let two boundaries be equal or
differ in one non-star position. Any overlapping component pair has
at least ten common vertices.
\end{lemma}
\begin{proof}
An unchanged constant component has 15 vertices. Between a constant
and nonconstant boundary, one additional colour is excluded at one
coordinate of that 15-cycle. A cyclic permutation of its three colours
preserves each component: an order-three action on a two-element
component set is trivial. Hence each colour appears five times at
each coordinate and precisely ten states remain.
In every other case, both components are their entire local state
sets. The intersection lists have size at least three except possibly
one list of size two. Start at that position, proceed around the
cycle, and finally avoid the two already coloured neighbours.
This gives at least $2\cdot2^3\cdot1=16$ colourings. With no changed
position the bound only improves. Full lists cause no difficulty.
\end{proof}

Consequently, if each of two boundaries changes at most one entry,
normalize the target \emph{relative to the source maps}. The two
reference product cycles have at least ten common internal edges:
one left protected edge from Lemma~\ref{lem:pentagon-protection},
copied at the ten or more common right states from
Lemma~\ref{lem:pentagon-overlap}.
The two coordinate permutations and colour permutations may be
independent. Actual synchronous changes are included in this stronger
local classification; physical exterior vertices are never identified
merely because they have equal colours.

\subsection{An arbitrary bottom cycle}
\begin{lemma}\label{lem:two-pentagons-restore}
Suppose $G$ has girth at least five and contains two vertex-disjoint
5-cycles with no edge between them. Each cycle has vertices of degree
three except for at most one vertex of degree two.
Let $H$ be obtained by deleting both cycles. If $\R_4(H)$ is
Hamiltonian, so is $\R_4(G)$.
Neither planarity nor three-colourability of $H$ is assumed here.
\end{lemma}
\begin{proof}
Each degree-three cycle vertex has one exterior neighbour; a
degree-two vertex gives a full list. An exterior vertex meets a
given 5-cycle at most once, since two attachments would create a
3- or 4-cycle. Thus a bottom recolouring changes at most one list
in each factor. The two factors may change simultaneously.
All exterior neighbours remain in $H$, which is nonempty.

Along a bottom Hamilton cycle, every ordinary pentagon has a
nonconstant boundary somewhere. Otherwise consecutive constant
five-tuples, differing in at most one entry, must be equal.
Its entire boundary would stay fixed along the bottom cycle.
A global colour permutation changing that colour supplies another
proper colouring of $H$ which the Hamilton cycle must visit,
a contradiction. These nonconstant layers need not coincide
for the two factors.

Make a component graph whose vertices are all product components
in all bottom layers, joining components in consecutive layers
when they have a common local state. There are at most four components
per layer and every component-graph degree is at most eight.
This graph is connected. Indeed, on each constant run the two
15-cycle components can be labelled consistently by a binary bit:
consecutive constant boundaries are equal. A nonconstant layer
forgets, or resets, that bit. A starred factor resets it at every layer.
Take four bit-labelled copies of the bottom cycle and allow a flip of
each bit at its reset layers. This graph is connected because each
bit has a reset somewhere. Contract the reset edges within each
layer. The result is exactly the component graph: intersections
between nonconstant components are nonempty, each constant component
meets its nonconstant neighbour, and an unchanged constant component
meets precisely itself.

Choose a spanning tree of the component graph. Give the root one
product cycle from Lemma~\ref{lem:protected-product} and record its
two restoration maps. For a child, apply the inverse parent maps
to the target boundaries, normalize, then compose restoration maps.
The target component is fixed by any common state. This makes the
source a standard type for the protected-edge table, so the
two \emph{reference} product cycles share at least ten internal edges.
The parent's component has tree degree at most eight and has so far
lost at most seven internal edges. The child is fresh.
Choose an unused common edge $xy$, delete its two copies and add
the cross-layer edges at $x$ and $y$. Each cross edge is the actual
bottom recolouring with the complete local colouring unchanged.
The two circles become two paths joined into one circle.
Induction over the tree preserves one circle, disjoint coverage and
legality, and ends with all extensions of all bottom colourings.
There is no additional unproved closing step.
\end{proof}

\subsection{The planar structural reduction}
\begin{lemma}\label{lem:two-separated-five-faces}
Every connected plane graph of girth at least five with
$2\le d(v)\le3$ and at most one degree-two vertex contains
two vertex-disjoint 5-cycles with no edge between them.
\end{lemma}
\begin{proof}
For any connected cyclic plane graph of girth at least five and
maximum degree three, let $p$ and $q$ count vertices of degrees one
and two. Euler's identity, $2e=3n-2p-q$, and the face-length sum give
\[
 f_5=12-4p-2q+\sum_{|F|\ge7}(|F|-6).                 \tag{*}
\]
Every facial walk contains a cycle and so has length at least five;
bridges count twice. A length-five facial walk is a simple 5-cycle.
For the graph in the lemma, $p=0,q\le1$, hence $f_5\ge12-2q$.

If there is a bridge $ab$, take the side $C$ not containing the
unique degree-two vertex, if it exists. In $C$, only $a$ has degree
two and all other vertices have degree three. This side is cyclic,
and (*) gives at least ten 5-faces. At most two contain $a$.
Choose a 5-cycle $S$ avoiding $a$. The other bridge side has either
at most two degree-two vertices, or one degree-one vertex when
$b$ was the original degree-two vertex. All other vertices have
degree three. The degree sum rules out a tree and (*) gives at
least eight 5-faces. Any such cycle $P$ is disjoint from $S$ and
has no edge to it, because $ab$ was the only edge between the sides.

Assume now there is no bridge. The subcubic graph is 2-connected:
at a cutvertex every branch would require at least two incident
edges to avoid a bridge, forcing degree at least four.
Take a facial 5-cycle $S$ and let $B$ be its exterior neighbours.
Girth makes them distinct; $b=|B|$ is five, or four when $S$
contains the unique degree-two vertex $z$.
Set $H=G-S$, embedded with the deleted region outside.
Every member of $B$ is on that outside face. Its degree drops
from three to two, except that $z\in B$ would become a leaf.
There is at most one leaf and no isolated vertex, so every component
of $H$ is cyclic.

In one such component, let $r$ count its members of $B$, $q'\le1$
its original degree-two vertices, and $\ell$ its outside facial
length. Formula (*) after removing the outside-face contribution is
\[
 f^{\rm in}_5=\ell+6-2r-2q'+L,\qquad
 L=\sum_{\substack{F\text{ internal}\\|F|\ge7}}(|F|-6)\ge0.       \tag{**}
\]
Even when $z\in B$, the degree deficit $4p+2q$ remains
$2r+2q'$, so the formula still applies. Here $\ell\ge5$.
A member of $B$ belongs to at most one internal 5-face: its degree
is at most two, with an outside face angle, and a leaf belongs to none.
If $H$ is disconnected, some component has $r\le2$ since
$|B|\le5$. By (**) it has at least $9-2r>r$ internal 5-faces.
One avoids $B$, and is the required $P$.

Suppose $H$ is connected. The $b$ exterior edges partition the
outside angles along $S$ into $b$ face portions, so at most $b$
other faces are involved. Deleting five vertices and $5+b$ edges
decreases the face count by $b$; thus these $b$ neighbours
$F_1,\ldots,F_b$ are distinct. Their merger has outside length
\[
 \ell=\sum_{i=1}^b|F_i|-(2b+5),                       \tag{***}
\]
because the five cycle edges are removed once and the $b$ outgoing
edges twice from that sum.

If $q=0$, any $S$ has $b=5$, $\ell\ge10$, $q'=0$.
Formula (**) gives at least six internal 5-faces, more than $|B|$.
If $z$ lies on a 5-face, choose that face as $S$. Now $b=4$,
$\ell\ge7$, $q'=0$, and there are at least five, again more than $|B|$.

It remains that $z$ lies on no 5-face. There is a 5-face adjacent
to two distinct non-5-faces. To see this, consider the plane dual.
It is loopless because the original graph is bridgeless.
All its faces are triangles except one digon corresponding to $z$.
The distinct endpoints of that digon both represent non-5-faces.
Let $W$ be the set of all such dual vertices. If every vertex outside
$W$ had at most one distinct neighbour in $W$, a triangular face
incident with a $W$--$W$ edge would have its third vertex in $W$.
Starting with a digon edge and going around either endpoint therefore
shows that all its neighbours are in $W$: a triangular face propagates
the property to the next neighbour, and the digon merely repeats the
same neighbour. Repeating at any reached neighbour makes $W$ closed
under adjacency. Dual connectivity would force all vertices into $W$,
contradicting $f_5\ge10$. This proves the assertion.
Choose the resulting 5-face $S$. If its deletion disconnects the graph,
use the case already proved; otherwise $b=5,q'=1$ and two distinct
neighbours in (***) have length at least six. Thus $\ell\ge12$ and
(**) again gives at least six internal 5-faces. In every case one
avoids $B$, making it disjoint and nonadjacent to $S$.
\end{proof}

\subsection{Induction, checks and the remaining frontier}
\begin{proof}[Proof of Theorem~\ref{thm:subcubic}]
Induct on the number of vertices, allowing disconnected graphs.
The one-vertex case is $\R_4(G)=K_4$. A vertex of degree at most one
is handled by Lemma~\ref{lem:single-lift}.
If there are two degree-two vertices, adjacent ones use Lemma~\ref{lem:adjacent-two-four}. Ones at distance two use Lemma~\ref{lem:path232}, unless their
middle vertex has degree two, when an adjacent pair already applies.
Ones at distance at least three use Lemma~\ref{lem:separated-two-g5}, whose spare-colour witness uses
Gr\"otzsch's theorem~\cite{Grotzsch59}; the new pentagon lift does not.
Pairs in different components have disjoint neighbourhoods and
are covered by the same separated-pair argument.

The remaining graph has minimum degree at least two and at most
one degree-two vertex. Apply Lemma~\ref{lem:two-separated-five-faces}
inside any component. Delete its two separated pentagons.
The remaining graph is nonempty and belongs to the same class, so
induction gives its four-colour Hamilton cycle.
Lemma~\ref{lem:two-pentagons-restore} restores both pentagons.
This proves the assertion at four colours.
For every $k\ge5$ use the triangle-free planar theorem proved earlier,
with larger colours separately supplied there by degeneracy.
The frozen $012012$ colouring of $C_6$ gives the Hamiltonian
threshold lower bound four. No monotonicity in $k$ is presumed.
\end{proof}



\subsection{Block-subcubic lemma}
\begin{lemma}\label{lem:block-subcubic}
Let $G$ be a nonempty finite simple planar graph of girth at least
five. If every block $B$ satisfies $\Delta(B)\le3$, then $\Rk(G)$ is
Hamiltonian for every $k\ge4$. Cutvertices may have arbitrarily large
degree in $G$. The uniform Hamiltonian threshold of this class is exactly four.
Four-colour Hamilton-connectedness is not asserted.
\end{lemma}
\begin{proof}
Treat bridges as two-vertex blocks and isolated vertices as singleton
blocks. The class is closed under vertex deletion: a nontrivial
two-connected block of a subgraph is a two-connected subgraph of
$G$, hence is contained in one original block, with no larger
internal degrees. Planarity and girth are also preserved.

Products of Hamiltonian graphs are Hamiltonian: copy one factor
cycle along a spanning path of another factor cycle, and merge
successive copies using two alternating, distinct internal edges.
This also handles disconnected original graphs. Choose a smallest
four-colour counterexample, which may therefore be assumed connected.
The singleton base is $K_4$, and the single-vertex lift excludes
vertices of degree at most one. If there is only one block, the
previous subcubic girth-five theorem applies. Thus assume minimum
degree at least two and at least one cutvertex.

Lemmas~\ref{lem:adjacent-two-four}, \ref{lem:path232}, and
\ref{lem:separated-two-g5} exclude adjacent degree-two vertices, a
$(2,3,2)$ path, and two degree-two vertices at distance at least three.
Consequently any two degree-two vertices not already giving a
reduction must have a common neighbour of degree at least four in $G$.

The block-cut tree has at least two endblocks. No endblock is a
bridge, since its noncut endpoint would be a leaf. Let $B$ be an
endblock and $x$ its unique cutvertex of $G$. Every vertex of
$B-x$ has its full $G$-degree in $B$, hence degree two or three;
also $d_B(x)\in\{2,3\}$. Let $q$ count the degree-two vertices
of $B-x$ and put $\varepsilon=[d_B(x)=2]$. Euler's identity in $B$
gives
\[
 f_5(B)=12-2(q+\varepsilon)+\sum_{|F|\ge7}(|F|-6).
\]
Each five-face is a simple pentagon, and at most $d_B(x)\le3$
five-faces meet $x$.

If $q\le1$, then $f_5(B)\ge8>3$, so choose a five-face avoiding
$x$. Its vertices have degree at most three in $G$, with at most
one degree-two vertex. If $q\ge2$, take any two degree-two vertices
of $B-x$. Their neighbours all lie in $B$. In the absence of the
three preceding reductions, a common neighbour has $G$-degree at
least four; the only possible such vertex in this endblock is $x$.
Thus all $q$ degree-two vertices are adjacent to $x$, and
$q\le d_B(x)\le3$. In either possible degree of $x$ this gives
$q+\varepsilon\le3$, whence $f_5(B)\ge6>3$. A five-face avoiding
$x$ also avoids all those degree-two vertices, because a cycle
through any one of them must use its edge to $x$.

Choose these pentagons in two distinct endblocks. Their vertices
are noncutvertices of $G$ and belong to their respective endblocks
only, so the pentagons are disjoint and have no edge between them.
The established protected-edge lift for two separated pentagons
applies, including its allowance of one degree-two vertex per
pentagon and simultaneous exterior list changes. The smaller graph
remains in the class and has a bottom Hamilton cycle by induction.
Restoration contradicts minimality.

For every $k\ge5$ separately use the established triangle-free
planar theorem. Finally $C_6$ belongs to the class and its periodic
$012012$ colouring is frozen in $\R_3(C_6)$, giving the lower bound
four. No monotonicity in the number of colours has been assumed.
\end{proof}
In particular, a remaining minimum counterexample for general
planar girth-five four-colour Hamiltonicity must contain a block
with a vertex of degree at least four \emph{inside that block}.
High total degree accumulated only at cutvertices is insufficient.
This does not assert that every minimum counterexample is two-connected.

\section{Sharpness and the unresolved frontier}
\label{sec:sharpness}

The lower bounds used above are frozen colourings, not claims about a particular embedding.

\begin{proposition}[Sharpness examples]
\label{prop:sharpness}
The following obstructions hold.
\begin{enumerate}[label=(\roman*)]
\item The icosahedral graph has a frozen proper six-colouring (see Bousquet--Perarnau~\cite{BP16}), so $\R_6$ is disconnected.  Hence the planar threshold is at least seven.
\item The cube has a frozen proper four-colouring (identify colours with $\mathbb Z_2^2$ and use the three nonzero edge directions), so triangle-free and bipartite planar classes have threshold at least five.
\item A cycle whose length is a multiple of three has a frozen periodic three-colouring.  This gives the lower bound four for girth at least six and for the four-colour girth-five subclasses containing $C_6$.
\item $\R_3(P_3)$ at $k=3$ has twelve states; its six degree-two states force three disjoint four-cycles, so it has no spanning cycle.  Also $\R_3(C_5)=C_{15}\sqcup C_{15}$.  Thus the forest and $C_5$ thresholds cannot be three.
\end{enumerate}
\end{proposition}

\begin{proof}
For (i)--(iii), every vertex sees all colours other than its own, so the displayed colouring is frozen.  For (iv), the forced-edge description of $\R_3(P_3)$ gives three disjoint four-cycles, while the thirty proper three-colourings of $C_5$ have degree two and split into two components of length fifteen. \qed
\end{proof}

The general planar girth-five problem is deliberately left at the following boundary:
\[
 4\le h(\mathcal P_5)\le5,\qquad
 4\le\hc(\mathcal P_5)\le5,
\]
where $\mathcal P_5$ is the class of planar graphs of girth at least five.  The upper bounds are the five-colour Hamilton-connected theorem.  The lower bounds record the four-colour boundary; no complete-state-graph counterexample has been verified in the present record.

The following implications are invalid and are not used here:
\begin{itemize}
\item a connected or 2-connected recolouring graph need not have a Hamilton cycle;
\item a Hamilton cycle in a fixed list fibre need not extend to the complete recolouring graph;
\item a spanning two-factor, a quotient cycle, or a finite local kernel does not by itself give a single Hamilton cycle;
\item the four-colour result for at most seven five-cycles or for subcubic graphs does not settle all planar graphs of girth at least five.
\end{itemize}

\section{Verification note}
\label{sec:verification}

The finite data used for the endpoint-sensitive statements were replayed in fresh output directories.  The local Hamilton-connected kernel for the triangle-free five-colour induction consists of 53 representatives, 225 labelled list pairs, and 3225 legal interfaces; the independent reader checked 27,600 entry cases.  These numbers certify the displayed finite local templates only.  The complete $C_5$ certificate contains 240 states, 478 prescribed-endpoint paths, and 114,242 single-coordinate steps; the two three-colour components were read back as $15$-cycles.  For the protected pentagon lift, the fresh readers covered 590 ordinary factor words, 465 factor cycles, 475 seeds, 1024 boundary records, and 245,760 labelled states; the protected reader covered 23 component types, 2304 labelled boundaries, and 313 interfaces, with at least ten common states per interface.  The block-structure reader checked 22 structures, including a candidate whole graph with 1,495,752 states.  These are bounded certificate replays, not arbitrary-size graph searches.

All arbitrary-order assertions in this paper are supplied by the analytic inductions: the degenerate theorem, the separated-degree arguments in the planar Hamilton-connected proofs, the Euler reductions for the few-five-cycle class, and the protected-edge restoration for the subcubic classes.  The finite checks do not constitute external peer review, formal verification, or an extrapolation from sampled graphs.  The current research record keeps the complete planar girth-five four-colour Hamiltonian and Hamilton-connected statements as unresolved.

\appendix
\section{Literal pentagon factors and common-port seeds}
\label{app:pentagon-factor-words}
Five digits record the five colours in cyclic vertex order. In six-digit
words the first digit is the path layer. Semicolons separate factor blocks;
without a semicolon the entire factor line is one block.
All Hamilton words are cyclic. A port is a pair of five-digit words.
These literal identities are part of the proof, not solver verdicts.
\paragraph{Factor 00000 (30 words).}
% BEGIN FACTOR-00000
\begingroup\footnotesize\ttfamily\begin{flushleft}
12123 12323 12313 12312 32312 31312 31212 31232 31231 21231 23231 23131 23121 23123 13123\\
12132 13132 13232 13212 13213 23213 21213 21313 21323 21321 31321 32321 32121 32131 32132\\
\end{flushleft}\endgroup
% END FACTOR-00000
\paragraph{Factor 00001 (44 words).}
% BEGIN FACTOR-00001
\begingroup\small\ttfamily\begin{flushleft}
12120 32120 ; 12123 12323 ; 12130 32130 ; 12132 32132 ; 12310 32310\\
12312 12313 ; 12320 32320 ; 13120 23120 ; 13123 23123 ; 13130 23130\\
13132 13232 ; 13210 23210 ; 13212 13213 ; 13230 23230 ; 21210 31210\\
21213 23213 ; 21230 31230 ; 21310 31310 ; 21313 21323 ; 21320 31320\\
31212 31232 ; 31312 32312\\
\end{flushleft}\endgroup
% END FACTOR-00001
\paragraph{Factor 00011 (42 words).}
% BEGIN FACTOR-00011
\begingroup\small\ttfamily\begin{flushleft}
12102 32102 ; 12103 12123 ; 12120 32120 ; 12130 32130 ; 12132 32132\\
12302 12303 ; 12320 12323 ; 13102 13202 ; 13103 23103 ; 13120 23120\\
13123 23123 ; 13130 23130 ; 13132 13232 ; 13203 23203 ; 13230 23230\\
21203 21303 ; 21230 31230 ; 21320 21323 ; 31202 31232 ; 31302 32302\\
31320 32320\\
\end{flushleft}\endgroup
% END FACTOR-00011
\paragraph{Factor 00012 (47 words).}
% BEGIN FACTOR-00012
\begingroup\small\ttfamily\begin{flushleft}
12103 12303 ; 12120 12320 ; 12123 12323 ; 12130 32130 ; 13103 13203\\
13120 13130 ; 13123 23123 ; 13230 23230 ; 21201 31201 ; 21203 23203\\
21230 31230 ; 21231 31231 ; 21301 21303 ; 23101 23103 ; 23120 23121\\
23130 23131 ; 23201 23231 ; 31301 32301 ; 31320 32320 ; 31321 32321\\
32101 32131 ; 32120 32121 ; 21320 21321 21323\\
\end{flushleft}\endgroup
% END FACTOR-00012
\paragraph{Factor 00101 (64 words).}
% BEGIN FACTOR-00101
\begingroup\small\ttfamily\begin{flushleft}
12010 32010 ; 12012 32012 ; 12013 12313 ; 12020 32020 ; 12023 12323\\
12030 32030 ; 12032 32032 ; 12310 32310 ; 12312 32312 ; 12320 32320\\
13010 23010 ; 13012 13212 ; 13013 23013 ; 13020 23020 ; 13023 23023\\
13030 23030 ; 13032 13232 ; 13210 23210 ; 13213 23213 ; 13230 23230\\
21010 31010 ; 21013 21313 ; 21020 31020 ; 21023 21323 ; 21030 31030\\
21210 21213 ; 21230 31230 ; 21310 31310 ; 21320 31320 ; 31012 31312\\
31032 31232 ; 31210 31212\\
\end{flushleft}\endgroup
% END FACTOR-00101
\paragraph{Factor 00102 (63 words).}
% BEGIN FACTOR-00102
\begingroup\small\ttfamily\begin{flushleft}
12013 12023 ; 12310 12313 ; 12320 12323 ; 13010 23010 ; 13013 23013\\
13020 23020 ; 13023 23023 ; 13030 13230 ; 13210 13213 ; 21010 21013\\
21020 21023 ; 21021 23021 ; 21030 23030 ; 21031 31031 ; 21210 31210\\
21213 21313 ; 21230 23230 ; 21231 31231 ; 21310 21320 ; 21321 21323\\
23031 23231 ; 23210 23213 ; 31010 32010 ; 31020 31320 ; 31021 32021\\
31030 31230 ; 31310 32310 ; 31321 32321 ; 32020 32320 ; 32030 32031\\
12010 12020 12030\\
\end{flushleft}\endgroup
% END FACTOR-00102
\paragraph{Factor 00112 (51 words).}
% BEGIN FACTOR-00112
\begingroup\small\ttfamily\begin{flushleft}
12020 12023 ; 12030 32030 ; 12303 12323 ; 12320 32320 ; 13020 13023\\
13030 13230 ; 13203 23203 ; 21030 31030 ; 21031 23031 ; 21201 23201\\
21203 21303 ; 21230 31230 ; 21231 31231 ; 21301 21321 ; 21320 21323\\
23020 23030 ; 23021 23023 ; 23230 23231 ; 31020 32020 ; 31021 32021\\
31031 32031 ; 31201 31301 ; 31320 31321 ; 32301 32321 ; 21020 21021 21023\\
\end{flushleft}\endgroup
% END FACTOR-00112
\paragraph{Factor 00121 (58 words).}
% BEGIN FACTOR-00121
\begingroup\small\ttfamily\begin{flushleft}
12010 32010 ; 12012 12013 ; 12030 32030 ; 12032 32032 ; 12302 32302\\
12303 12313 ; 12310 32310 ; 12312 32312 ; 13010 23010 ; 13012 13212\\
13013 13213 ; 13030 23030 ; 13032 13232 ; 13202 13203 ; 13210 23210\\
13230 23230 ; 21010 31010 ; 21013 23013 ; 21030 31030 ; 21203 21303\\
21210 21213 ; 21230 31230 ; 21310 21313 ; 23203 23213 ; 31012 32012\\
31032 31232 ; 31202 31302 ; 31210 31212 ; 31310 31312\\
\end{flushleft}\endgroup
% END FACTOR-00121
\paragraph{Factor 00123 (57 words).}
% BEGIN FACTOR-00123
\begingroup\small\ttfamily\begin{flushleft}
12010 12310 ; 12012 32012 ; 12030 13030 ; 12032 32032 ; 12302 12312\\
13010 23010 ; 13012 13032 ; 13210 13230 ; 21010 31010 ; 21030 31030\\
21031 23031 ; 21201 31201 ; 21210 23210 ; 21230 31230 ; 21231 31231\\
21301 31301 ; 21310 31310 ; 23030 23230 ; 23201 23231 ; 31012 31312\\
31031 32031 ; 31032 31232 ; 31202 31302 ; 31210 31212 ; 32010 32030\\
32301 32302 ; 32310 32312 ; 13202 13212 13232\\
\end{flushleft}\endgroup
% END FACTOR-00123
\paragraph{Factor 01012 (68 words).}
% BEGIN FACTOR-01012
\begingroup\small\ttfamily\begin{flushleft}
10103 20103 ; 10120 30120 ; 10123 20123 ; 10130 30130 ; 10203 20203\\
10230 30230 ; 10303 20303 ; 10320 30320 ; 10323 20323 ; 12103 12303\\
12120 32120 ; 12123 12323 ; 12130 32130 ; 12320 32320 ; 13103 23103\\
13120 23120 ; 13123 23123 ; 13130 23130 ; 13203 23203 ; 13230 23230\\
20101 20131 ; 20120 20320 ; 20121 23121 ; 20130 20230 ; 20201 30201\\
20231 30231 ; 20301 20321 ; 23101 23201 ; 23131 23231 ; 30101 32101\\
30121 32121 ; 30131 32131 ; 30301 32301 ; 30321 32321\\
\end{flushleft}\endgroup
% END FACTOR-01012
\paragraph{Factor 01023 (66 words).}
% BEGIN FACTOR-01023
\begingroup\small\ttfamily\begin{flushleft}
10102 30102 ; 10130 30130 ; 10132 30132 ; 10202 30202 ; 10210 30210\\
10212 30212 ; 10230 30230 ; 10232 30232 ; 10302 30302 ; 10310 30310\\
10312 30312 ; 12102 32102 ; 12130 32130 ; 12132 13132 ; 12302 32302\\
12310 32310 ; 12312 32312 ; 13102 13202 ; 13130 23130 ; 13210 23210\\
13212 13232 ; 13230 23230 ; 20101 20301 ; 20130 20230 ; 20131 30131\\
20201 30201 ; 20210 20310 ; 20231 30231 ; 23101 23201 ; 23131 23231\\
30101 32101 ; 30301 32301 ; 32131 32132\\
\end{flushleft}\endgroup
% END FACTOR-01023
\paragraph{Cycle 00011 (42 words).}
% BEGIN C1-00011
\begingroup\small\ttfamily\begin{flushleft}
12102 13102 13103 23103 23123 13123 13120 23120 23130 13130\\
13230 23230 21230 31230 31232 31202 31302 32302 12302 12303\\
12103 12123 12323 12320 32320 31320 21320 21323 21303 21203\\
23203 13203 13202 13232 13132 12132 12130 12120 32120 32130\\
32132 32102\\
\end{flushleft}\endgroup
% END C1-00011
\paragraph{Cycle 00012 (47 words).}
% BEGIN C1-00012
\begingroup\small\ttfamily\begin{flushleft}
12103 12123 13123 13120 12120 32120 32130 12130 13130 13230\\
23230 21230 31230 31231 21231 23231 23201 23101 23131 23130\\
23120 23121 23123 23103 13103 13203 23203 21203 21201 31201\\
31301 32301 32101 32131 32121 32321 31321 21321 21301 21303\\
21323 21320 31320 32320 12320 12323 12303\\
\end{flushleft}\endgroup
% END C1-00012
\paragraph{Cycle 00101 (64 words).}
% BEGIN C1-00101
\begingroup\small\ttfamily\begin{flushleft}
12010 12020 32020 31020 21020 23020 13020 13010 23010 23210\\
13210 13212 13232 13032 13012 12012 32012 31012 31312 32312\\
12312 12310 32310 31310 21310 21320 31320 32320 12320 12323\\
12313 12013 12023 13023 23023 21023 21323 21313 21013 23013\\
13013 13213 23213 21213 21210 31210 31212 31232 31032 32032\\
12032 12030 32030 31030 31230 21230 23230 13230 13030 23030\\
21030 21010 31010 32010\\
\end{flushleft}\endgroup
% END C1-00101
\paragraph{Cycle 00102 (63 words).}
% BEGIN C1-00102
\begingroup\small\ttfamily\begin{flushleft}
12010 12030 32030 31030 31230 21230 23230 13230 13030 23030\\
21030 21031 23031 23231 21231 31231 31031 32031 32021 32321\\
31321 21321 21320 31320 32320 12320 12310 32310 31310 21310\\
21313 21323 21023 23023 13023 13020 13010 23010 21010 21020\\
23020 23021 21021 31021 31020 32020 12020 12023 12323 12313\\
12013 13013 23013 21013 21213 23213 13213 13210 23210 21210\\
31210 31010 32010\\
\end{flushleft}\endgroup
% END C1-00102
\paragraph{Cycle 00121 (58 words).}
% BEGIN C1-00121
\begingroup\small\ttfamily\begin{flushleft}
12010 13010 23010 23210 13210 13230 13030 12030 32030 31030\\
31032 32032 12032 13032 13232 13202 13212 13012 12012 32012\\
31012 31312 32312 12312 12310 32310 31310 21310 21210 21213\\
23213 13213 13203 23203 21203 21303 21313 21013 23013 13013\\
12013 12313 12303 12302 32302 31302 31202 31232 31212 31210\\
31230 21230 23230 23030 21030 21010 31010 32010\\
\end{flushleft}\endgroup
% END C1-00121
\paragraph{Cycle 00123 (57 words).}
% BEGIN C1-00123
\begingroup\small\ttfamily\begin{flushleft}
12010 12310 32310 31310 21310 21010 23010 13010 13210 13212\\
13202 13232 13032 13012 12012 32012 31012 31212 31232 31032\\
32032 12032 12030 32030 32031 31031 21031 21030 31030 31230\\
21230 21231 31231 31201 31202 31302 31312 32312 12312 12302\\
32302 32301 31301 21301 21201 23201 23231 23031 23030 13030\\
13230 23230 23210 21210 31210 31010 32010\\
\end{flushleft}\endgroup
% END C1-00123
\paragraph{Cycle 01012 (68 words).}
% BEGIN C1-01012
\begingroup\small\ttfamily\begin{flushleft}
10103 10203 13203 23203 20203 20303 10303 12303 12103 12123\\
10123 13123 23123 20123 20323 10323 12323 12320 32320 30320\\
10320 20320 20120 23120 13120 10120 30120 32120 12120 12130\\
32130 30130 10130 10230 30230 20230 23230 13230 13130 23130\\
20130 20131 23131 23231 20231 30231 30131 32131 32101 32301\\
30301 20301 20321 30321 32321 32121 30121 20121 23121 23101\\
23201 20201 30201 30101 20101 20103 23103 13103\\
\end{flushleft}\endgroup
% END C1-01012
\paragraph{Cycle 01023 (66 words).}
% BEGIN C1-01023
\begingroup\small\ttfamily\begin{flushleft}
10102 10302 30302 32302 12302 12312 10312 30312 32312 32310\\
12310 10310 30310 20310 20210 30210 10210 10212 30212 30202\\
10202 13202 13102 13132 13232 13212 13210 23210 23230 23130\\
23131 23231 23201 23101 20101 30101 30201 20201 20301 30301\\
32301 32101 32131 30131 30231 20231 20131 20130 20230 30230\\
10230 13230 13130 12130 32130 30130 10130 10132 10232 30232\\
30132 32132 12132 12102 32102 30102\\
\end{flushleft}\endgroup
% END C1-01023
\paragraph{C2 for 00001 (88 words).}
% BEGIN C2-00001
\begingroup\small\ttfamily\begin{flushleft}
012120 012130 012132 032132 032130 032120 132120 132130 132132\\
112132 112130 112120 112123 012123 012323 012320 112320 112323\\
112313 112312 012312 012313 012310 112310 132310 032310 031310\\
021310 121310 131310 131312 132312 032312 031312 031212 031210\\
031230 031232 131232 131212 131210 131230 121230 121210 121213\\
121313 121323 121320 131320 132320 032320 031320 021320 021323\\
021313 021213 021210 021230 023230 123230 113230 113210 123210\\
123213 113213 113212 113232 113132 013132 013232 013212 013213\\
023213 023210 013210 013230 013130 023130 023120 023123 123123\\
123120 123130 113130 113120 113123 013123 013120\\
\end{flushleft}\endgroup
% END C2-00001
\paragraph{C3 for 00001 (132 words).}
% BEGIN C3-00001
\begingroup\small\ttfamily\begin{flushleft}
012120 012130 012132 032132 032130 032120 032320 031320 021320\\
021323 021313 021213 023213 023210 013210 013230 023230 021230\\
121230 121210 121213 121313 121323 121320 131320 132320 132120\\
132130 132132 112132 112130 112120 112320 112323 112313 112310\\
132310 132312 112312 212312 212313 212310 212320 212323 212123\\
112123 012123 012323 012320 012310 012313 012312 032312 032310\\
031310 131310 231310 232310 232312 231312 231212 231210 231230\\
231232 131232 031232 031230 131230 131210 131212 131312 031312\\
031212 031210 021210 021310 121310 221310 221210 221230 223230\\
123230 113230 213230 213130 223130 223120 223123 123123 123120\\
123130 113130 113120 113123 213123 213120 212120 212130 212132\\
232132 232130 232120 232320 231320 221320 221323 221313 221213\\
223213 223210 213210 213213 213212 213232 213132 113132 113232\\
113212 113210 123210 123213 113213 013213 013212 013232 013132\\
013130 023130 023120 023123 013123 013120\\
\end{flushleft}\endgroup
% END C3-00001
\paragraph{Common ports 00001 (20 edges).}
% BEGIN PORTS-00001
\begingroup\small\ttfamily\begin{flushleft}
12120-12130 12130-12132 12132-32132 12320-12323 13120-13123\\
13130-23130 13132-13232 13210-23210 13212-13232 21213-21313\\
21313-21323 21320-21323 21320-31320 23120-23123 23120-23130\\
23210-23213 31210-31212 31320-32320 32120-32130 32130-32132\\
\end{flushleft}\endgroup
% END PORTS-00001
\paragraph{C2 for 00112 (102 words).}
% BEGIN C2-00112
\begingroup\small\ttfamily\begin{flushleft}
012020 012030 032030 032031 032021 031021 021021 021031 031031\\
031231 031201 031301 032301 032321 032320 012320 012323 012303\\
112303 112323 112320 132320 132321 132301 131301 131201 131231\\
131230 031230 031030 021030 021230 121230 121030 131030 131020\\
121020 123020 123023 123021 123031 123030 123230 113230 113030\\
113020 113023 112023 012023 013023 013020 013030 013230 023230\\
023030 023031 023231 123231 121231 021231 021201 021301 021321\\
031321 131321 131021 131031 121031 121021 121023 121323 121303\\
021303 021203 121203 123203 113203 013203 023203 023201 123201\\
121201 121301 121321 121320 131320 031320 021320 021323 021023\\
023023 023021 023020 021020 031020 032020 132020 132021 132031\\
132030 112030 112020\\
\end{flushleft}\endgroup
% END C2-00112
\paragraph{C3 for 00112 (153 words).}
% BEGIN C3-00112
\begingroup\small\ttfamily\begin{flushleft}
012020 012023 013023 013020 013030 013230 023230 021230 021030\\
031030 032030 032031 032021 132021 132031 132030 131030 121030\\
121230 123230 113230 113030 113020 113023 112023 112323 112320\\
132320 132321 132301 032301 032321 032320 012320 012323 012303\\
112303 212303 212323 212320 232320 232020 231020 221020 223020\\
223023 221023 221323 221320 231320 131320 031320 021320 121320\\
121323 121321 121301 221301 221321 231321 231021 231031 231231\\
231201 231301 232301 232321 232021 232031 232030 212030 212020\\
212023 213023 213020 213030 213230 223230 223030 223031 221031\\
221021 223021 123021 123023 023023 023021 023031 123031 123231\\
023231 023201 123201 123203 223203 213203 113203 013203 023203\\
021203 121203 121201 021201 031201 031301 021301 021321 031321\\
131321 131301 131201 131231 121231 021231 031231 031230 131230\\
231230 231030 221030 221230 221231 223231 223201 221201 221203\\
221303 121303 021303 021323 021023 121023 121021 121031 021031\\
021021 031021 031031 131031 131021 131020 121020 123020 123030\\
023030 023020 021020 031020 032020 132020 112020 112030 012030\\
\end{flushleft}\endgroup
% END C3-00112
\paragraph{Common ports 00112 (18 edges).}
% BEGIN PORTS-00112
\begingroup\small\ttfamily\begin{flushleft}
12020-12030 12023-13023 12320-12323 12320-32320 13020-13023\\
13020-13030 13030-13230 13230-23230 21020-23020 21020-31020\\
21021-21031 21030-21230 21030-31030 21301-21321 31201-31301\\
32021-32031 32030-32031 32301-32321\\
\end{flushleft}\endgroup
% END PORTS-00112


\section{Literal protected-pentagon certificates}
\label{app:protected-pentagon-words}
Boundary digit 4 denotes the full list. Colour words use only 0,1,2,3.
In six-digit cycle words the first digit is the path layer.
Pool order matters: the first two edges (one for the hard types) join
path layers or seed blocks; the last six join the factor tree.
The protected set is the specified cycle edge set, or COMMON set
for a hard type, minus POOL. The earlier appendix supplies the
ordinary cycles, the two constant-type 15-cycles, and the 00112 seeds.
An interface record gives source type:component, position, new colour,
normalized target type:component, and the actual common edge.
Positions start at zero. Two records per line are separated by semicolons.
\paragraph{PC2 for 00001 (88 words).}
% BEGIN PC2-00001
\begingroup\footnotesize\ttfamily\begin{flushleft}
012120 012130 112130 112120 132120 132130 132132 112132 113132 013132\\
012132 032132 032130 032120 032320 032310 012310 012313 012323 012320\\
112320 132320 131320 031320 021320 121320 121310 021310 021313 021323\\
121323 121313 121213 121210 131210 131212 131232 131230 121230 123230\\
113230 113232 113212 113210 123210 123213 113213 013213 023213 021213\\
021210 023210 013210 013212 013232 013230 023230 021230 031230 031232\\
031212 031210 031310 031312 032312 012312 112312 132312 131312 131310\\
132310 112310 112313 112323 112123 012123 013123 023123 023120 023130\\
013130 113130 123130 123120 123123 113123 113120 013120\\
\end{flushleft}\endgroup
% END PC2-00001
\paragraph{PC3 for 00001 (132 words).}
% BEGIN PC3-00001
\begingroup\footnotesize\ttfamily\begin{flushleft}
012120 013120 023120 023130 013130 012130 112130 113130 113132 112132\\
132132 132130 132120 232120 232130 232132 212132 213132 213232 213230\\
223230 221230 231230 231232 231212 231210 231310 231312 232312 232310\\
212310 212312 212313 212323 212123 112123 112120 113120 213120 212120\\
212130 213130 223130 123130 123120 223120 223123 213123 113123 123123\\
023123 013123 012123 012323 012313 012312 112312 132312 032312 032310\\
012310 012320 032320 132320 132310 112310 112313 112323 112320 212320\\
232320 231320 221320 221323 221313 221310 221210 221213 223213 123213\\
023213 021213 021210 023210 013210 013212 013213 113213 213213 213212\\
213210 223210 123210 113210 113212 113232 113230 123230 121230 131230\\
131232 131212 131210 121210 121213 121313 121323 021323 021313 021310\\
121310 121320 021320 031320 131320 131310 131312 031312 031310 031210\\
031212 031232 031230 021230 023230 013230 013232 013132 012132 032132\\
032130 032120\\
\end{flushleft}\endgroup
% END PC3-00001
\paragraph{Balanced common edges for 00001.}
% BEGIN COMMON-00001
\begingroup\footnotesize\ttfamily\begin{flushleft}
12132-13132 12132-32132 12310-32310 12313-12323 13123-23123\\
13210-13212 13210-23210 13230-13232 13230-23230 21210-21213\\
21230-23230 21230-31230 21313-21323 31210-31212 31212-31232\\
31230-31232 31310-31312 32120-32130 32130-32132\\
\end{flushleft}\endgroup
% END COMMON-00001
\paragraph{Cycle for 00004 (54 words).}
% BEGIN FHC-00004
\begingroup\footnotesize\ttfamily\begin{flushleft}
12120 12130 12132 13132 13130 23130 23230 23210 13210 13230\\
13232 13212 13213 23213 21213 21210 31210 31212 31232 31231\\
31230 21230 21231 23231 23131 23121 23123 23120 13120 13123\\
12123 12323 12313 12312 12310 32310 32312 31312 31310 21310\\
21313 21323 21321 21320 31320 31321 32321 32121 32131 32132\\
32130 32120 32320 12320\\
\end{flushleft}\endgroup
% END FHC-00004
\paragraph{Cycle for 00014 (60 words).}
% BEGIN FHC-00014
\begingroup\footnotesize\ttfamily\begin{flushleft}
12102 13102 13132 13232 13202 13203 13103 12103 12123 13123\\
13120 12120 32120 32320 12320 12323 12303 12302 32302 31302\\
31202 31232 31230 21230 23230 13230 13130 23130 23120 23121\\
23123 23103 23203 23201 23101 23131 23231 21231 31231 31201\\
21201 21203 21303 21323 21320 31320 31321 21321 21301 31301\\
32301 32321 32121 32101 32131 32130 12130 12132 32132 32102\\
\end{flushleft}\endgroup
% END FHC-00014
\paragraph{Cycle for 00104 (78 words).}
% BEGIN FHC-00104
\begingroup\footnotesize\ttfamily\begin{flushleft}
12010 12030 32030 31030 31230 21230 23230 13230 13030 23030\\
21030 21031 23031 23231 21231 31231 31031 32031 32021 32321\\
32320 12320 12323 12313 12013 12023 13023 23023 21023 21020\\
31020 32020 12020 13020 23020 23021 21021 31021 31321 31320\\
21320 21321 21323 21313 21310 31310 32310 12310 12312 32312\\
31312 31012 32012 12012 13012 13013 23013 21013 21213 23213\\
13213 13212 13232 13032 12032 32032 31032 31232 31212 31210\\
21210 23210 13210 13010 23010 21010 31010 32010\\
\end{flushleft}\endgroup
% END FHC-00104
\paragraph{Cycle for 00114 (62 words).}
% BEGIN FHC-00114
\begingroup\footnotesize\ttfamily\begin{flushleft}
12020 12320 32320 31320 21320 21323 21023 23023 23021 23031\\
21031 21021 21020 21030 21230 31230 31030 32030 12030 13030\\
13230 23230 23030 23020 13020 13023 12023 12323 12303 12302\\
32302 31302 31202 31232 31032 32032 12032 13032 13232 13202\\
13203 23203 21203 21303 21301 21321 31321 32321 32301 31301\\
31201 21201 23201 23231 21231 31231 31031 32031 32021 31021\\
31020 32020\\
\end{flushleft}\endgroup
% END FHC-00114
\paragraph{Cycle for 00124 (71 words).}
% BEGIN FHC-00124
\begingroup\footnotesize\ttfamily\begin{flushleft}
12010 12030 32030 32031 31031 31231 21231 21031 21030 31030\\
31230 21230 23230 13230 13030 23030 23031 23231 23201 23203\\
21203 21303 21301 21201 31201 31202 31302 31301 32301 32302\\
12302 12303 12313 12013 13013 13010 13210 13212 13012 12012\\
32012 31012 31312 32312 12312 12310 32310 31310 31210 31212\\
31232 31032 32032 12032 13032 13232 13202 13203 13213 23213\\
21213 21313 21013 23013 23010 23210 21210 21310 21010 31010\\
32010\\
\end{flushleft}\endgroup
% END FHC-00124
\paragraph{Cycle for 01014 (88 words).}
% BEGIN FHC-01014
\begingroup\footnotesize\ttfamily\begin{flushleft}
10102 10132 13132 13232 10232 30232 30132 32132 12132 12130\\
32130 30130 10130 10230 30230 20230 23230 13230 13130 23130\\
20130 20131 23131 23231 20231 30231 30131 32131 32101 32301\\
30301 20301 20101 30101 30201 20201 23201 23101 23121 20121\\
30121 32121 32321 30321 20321 20323 10323 12323 12320 32320\\
30320 10320 20320 20120 23120 13120 12120 32120 30120 10120\\
10123 12123 13123 23123 20123 20103 20203 20303 10303 10302\\
30302 32302 12302 12303 12103 10103 10203 13203 23203 23103\\
13103 13102 13202 10202 30202 30102 32102 12102\\
\end{flushleft}\endgroup
% END FHC-01014
\paragraph{Cycle for 01024 (85 words).}
% BEGIN FHC-01024
\begingroup\footnotesize\ttfamily\begin{flushleft}
10102 10202 30202 30212 10212 13212 13210 23210 20210 30210\\
10210 10213 10203 20203 23203 13203 13213 23213 20213 20313\\
10313 12313 12312 12310 32310 30310 20310 10310 10312 30312\\
32312 32302 30302 10302 10303 20303 20103 23103 13103 10103\\
12103 12303 12302 12102 32102 30102 30101 30201 20201 23201\\
23101 20101 20301 30301 32301 32101 32131 30131 30231 20231\\
23231 23131 20131 20130 23130 13130 13230 23230 20230 30230\\
10230 10130 30130 32130 12130 12132 32132 30132 30232 10232\\
10132 13132 13232 13202 13102\\
\end{flushleft}\endgroup
% END FHC-01024
\paragraph{Cycle for 01104 (74 words).}
% BEGIN FHC-01104
\begingroup\footnotesize\ttfamily\begin{flushleft}
10210 10310 20310 30310 32310 12310 12312 10312 30312 32312\\
32012 32032 32031 32021 32321 30321 20321 20320 10320 30320\\
32320 12320 12020 12010 32010 32020 32030 12030 12032 13032\\
13232 13212 13012 12012 12013 12313 10313 20313 20323 10323\\
12323 12023 13023 23023 23013 13013 13213 23213 20213 10213\\
10212 10232 10230 13230 23230 23030 13030 13020 23020 23021\\
23031 23231 20231 30231 30232 30212 30210 30230 20230 20210\\
23210 23010 13010 13210\\
\end{flushleft}\endgroup
% END FHC-01104
\paragraph{Cycle for 01124 (74 words).}
% BEGIN FHC-01124
\begingroup\footnotesize\ttfamily\begin{flushleft}
10202 10203 20203 23203 23201 20201 20301 20303 10303 12303\\
12302 32302 32301 30301 30201 30231 20231 20230 30230 10230\\
13230 23230 23231 23031 23030 13030 12030 32030 32031 32032\\
32012 32010 12010 12012 12032 13032 13232 10232 30232 30202\\
30212 10212 10213 10210 30210 20210 20310 30310 10310 12310\\
32310 32312 12312 12313 12013 13013 23013 23010 23210 13210\\
13010 13012 13212 13202 13203 13213 23213 20213 20313 10313\\
10312 30312 30302 10302\\
\end{flushleft}\endgroup
% END FHC-01124
\paragraph{Cycle for 01204 (83 words).}
% BEGIN FHC-01204
\begingroup\footnotesize\ttfamily\begin{flushleft}
10120 20120 23120 13120 13130 23130 20130 10130 30130 32130\\
12130 12132 32132 30132 10132 13132 13032 12032 32032 32012\\
12012 13012 13010 12010 32010 32020 12020 13020 23020 23010\\
23030 13030 12030 32030 32031 32131 30131 20131 23131 23031\\
23021 23121 20121 30121 32121 32021 32321 30321 20321 20320\\
10320 30320 32320 12320 12323 12023 13023 23023 23013 13013\\
12013 12313 12310 32310 32312 12312 10312 30312 30310 10310\\
20310 20313 10313 10323 20323 20123 23123 13123 10123 12123\\
12120 32120 30120\\
\end{flushleft}\endgroup
% END FHC-01204
\paragraph{Cycle for 01234 (82 words).}
% BEGIN FHC-01234
\begingroup\footnotesize\ttfamily\begin{flushleft}
10102 10302 30302 32302 12302 12303 10303 10323 20323 20303\\
20103 10103 12103 12123 13123 23123 20123 10123 10120 30120\\
32120 12120 13120 23120 20120 20320 10320 30320 32320 12320\\
12310 32310 30310 10310 20310 20313 10313 10312 30312 32312\\
12312 12313 12323 12023 13023 23023 23013 23010 13010 12010\\
12013 13013 13012 12012 32012 32010 32020 12020 13020 23020\\
23021 23121 20121 30121 32121 32021 32321 30321 20321 20301\\
30301 32301 32101 30101 20101 23101 23103 13103 13102 12102\\
32102 30102\\
\end{flushleft}\endgroup
% END FHC-01234
\paragraph{Pool for 00000:0.}
% BEGIN POOL-00000-0
\begingroup\footnotesize\ttfamily\begin{flushleft}
12123-12323 12123-13123 12312-12313 12313-12323\\
13123-23123 21231-31231 23121-23131 23131-23231\\
\end{flushleft}\endgroup
% END POOL-00000-0
\paragraph{Pool for 00000:1.}
% BEGIN POOL-00000-1
\begingroup\footnotesize\ttfamily\begin{flushleft}
13132-13232 13212-13213 21213-21313 21213-23213\\
21313-21323 21321-31321 32121-32131 32121-32321\\
\end{flushleft}\endgroup
% END POOL-00000-1
\paragraph{Pool for 00001:0.}
% BEGIN POOL-00001-0
\begingroup\footnotesize\ttfamily\begin{flushleft}
13210-13212 21210-21213 21230-31230 31210-31212\\
31230-31232 31310-31312 32120-32130\\
\end{flushleft}\endgroup
% END POOL-00001-0
\paragraph{Pool for 00011:0.}
% BEGIN POOL-00011-0
\begingroup\footnotesize\ttfamily\begin{flushleft}
12102-13102 13102-13103 13123-23123 13130-13230\\
13202-13203 21203-21303 31202-31302 32130-32132\\
\end{flushleft}\endgroup
% END POOL-00011-0
\paragraph{Pool for 00012:0.}
% BEGIN POOL-00012-0
\begingroup\footnotesize\ttfamily\begin{flushleft}
13130-13230 21301-21321 23120-23130 23121-23123\\
31201-31301 32101-32301 32121-32131 32121-32321\\
\end{flushleft}\endgroup
% END POOL-00012-0
\paragraph{Pool for 00101:0.}
% BEGIN POOL-00101-0
\begingroup\footnotesize\ttfamily\begin{flushleft}
12010-12020 12012-13012 12013-12023 12320-12323\\
13013-23013 13210-13212 21010-31010 23010-23210\\
\end{flushleft}\endgroup
% END POOL-00101-0
\paragraph{Pool for 00102:0.}
% BEGIN POOL-00102-0
\begingroup\footnotesize\ttfamily\begin{flushleft}
12013-13013 12313-12323 13013-23013 21013-21213\\
21023-21323 21031-23031 23020-23021 31031-31231\\
\end{flushleft}\endgroup
% END POOL-00102-0
\paragraph{Pool for 00112:0.}
% BEGIN POOL-00112-0
\begingroup\footnotesize\ttfamily\begin{flushleft}
12020-12030 12320-12323 13020-13030 21020-23020\\
21021-21031 21030-31030 32021-32031\\
\end{flushleft}\endgroup
% END POOL-00112-0
\paragraph{Pool for 00121:0.}
% BEGIN POOL-00121-0
\begingroup\footnotesize\ttfamily\begin{flushleft}
12010-13010 21010-21030 21010-31010 21210-21213\\
23010-23210 23030-23230 31012-31312 31030-32030\\
\end{flushleft}\endgroup
% END POOL-00121-0
\paragraph{Pool for 00123:0.}
% BEGIN POOL-00123-0
\begingroup\footnotesize\ttfamily\begin{flushleft}
13030-23030 13230-23230 21010-21310 21030-31030\\
21230-21231 31010-31210 31012-32012 31030-31230\\
\end{flushleft}\endgroup
% END POOL-00123-0
\paragraph{Pool for 01012:0.}
% BEGIN POOL-01012-0
\begingroup\footnotesize\ttfamily\begin{flushleft}
10103-13103 10123-13123 10130-30130 13130-13230\\
20101-30101 20230-30230 20301-20321 30131-30231\\
\end{flushleft}\endgroup
% END POOL-01012-0
\paragraph{Pool for 01023:0.}
% BEGIN POOL-01023-0
\begingroup\footnotesize\ttfamily\begin{flushleft}
12130-13130 20101-30101 20201-20301 20210-30210\\
20231-30231 23130-23131 30102-32102 30202-30212\\
\end{flushleft}\endgroup
% END POOL-01023-0
\paragraph{Pool for 00004:0.}
% BEGIN POOL-00004-0
\begingroup\footnotesize\ttfamily\begin{flushleft}
12312-12313 13130-13132 13212-13213 21210-21213\\
21321-21323 23121-23123 31231-31232 31310-31312\\
\end{flushleft}\endgroup
% END POOL-00004-0
\paragraph{Pool for 00014:0.}
% BEGIN POOL-00014-0
\begingroup\footnotesize\ttfamily\begin{flushleft}
12102-32102 12103-13103 12120-32120 12130-12132\\
13130-23130 21201-21203 23121-23123 32101-32121\\
\end{flushleft}\endgroup
% END POOL-00014-0
\paragraph{Pool for 00104:0.}
% BEGIN POOL-00104-0
\begingroup\footnotesize\ttfamily\begin{flushleft}
13012-13013 13210-23210 13212-13213 21020-21023\\
21320-21321 21321-21323 31320-31321 32320-32321\\
\end{flushleft}\endgroup
% END POOL-00104-0
\paragraph{Pool for 00114:0.}
% BEGIN POOL-00114-0
\begingroup\footnotesize\ttfamily\begin{flushleft}
12030-13030 21020-21021 21301-21321 21320-21323\\
23021-23023 31020-31021 31020-32020 32021-32031\\
\end{flushleft}\endgroup
% END POOL-00114-0
\paragraph{Pool for 00124:0.}
% BEGIN POOL-00124-0
\begingroup\footnotesize\ttfamily\begin{flushleft}
12310-32310 13210-13212 21030-31030 21201-31201\\
21213-21313 23010-23210 31010-32010 31301-31302\\
\end{flushleft}\endgroup
% END POOL-00124-0
\paragraph{Pool for 01014:0.}
% BEGIN POOL-01014-0
\begingroup\footnotesize\ttfamily\begin{flushleft}
10120-10123 10202-13202 10302-10303 12120-13120\\
12123-13123 13102-13103 20130-20131 20321-20323\\
\end{flushleft}\endgroup
% END POOL-01014-0
\paragraph{Pool for 01024:0.}
% BEGIN POOL-01024-0
\begingroup\footnotesize\ttfamily\begin{flushleft}
10103-13103 10132-13132 10302-10303 10312-30312\\
12130-12132 12312-12313 20103-20303 20213-20313\\
\end{flushleft}\endgroup
% END POOL-01024-0
\paragraph{Pool for 01104:0.}
% BEGIN POOL-01104-0
\begingroup\footnotesize\ttfamily\begin{flushleft}
10210-10310 12010-12020 12012-12013 12310-12312\\
13013-13213 20210-23210 23030-23230 32012-32312\\
\end{flushleft}\endgroup
% END POOL-01104-0
\paragraph{Pool for 01124:0.}
% BEGIN POOL-01124-0
\begingroup\footnotesize\ttfamily\begin{flushleft}
10202-10203 10210-10213 10230-13230 10310-30310\\
13010-13210 20210-20310 23230-23231 32310-32312\\
\end{flushleft}\endgroup
% END POOL-01124-0
\paragraph{Pool for 01204:0.}
% BEGIN POOL-01204-0
\begingroup\footnotesize\ttfamily\begin{flushleft}
10123-13123 12010-32010 12120-12123 12310-12313\\
12320-12323 23010-23020 30310-30312 32030-32031\\
\end{flushleft}\endgroup
% END POOL-01204-0
\paragraph{Pool for 01234:0.}
% BEGIN POOL-01234-0
\begingroup\footnotesize\ttfamily\begin{flushleft}
10102-30102 12102-32102 12123-13123 12302-32302\\
12310-12320 13012-13013 13102-13103 20101-30101\\
\end{flushleft}\endgroup
% END POOL-01234-0
\paragraph{All 313 single-pentagon interfaces.}
% BEGIN INTERFACES-ALL
\begingroup\footnotesize\ttfamily\begin{flushleft}
00000:0 0 1 00001:0 23121-23123 ; 00000:0 0 2 00001:0 31231-31232\\
00000:0 0 3 00001:0 21231-23231 ; 00000:0 1 1 00001:0 12312-32312\\
00000:0 1 2 00001:0 31212-31312 ; 00000:0 1 3 00001:0 31231-31232\\
00000:0 2 1 00001:0 21231-23231 ; 00000:0 2 2 00001:0 31312-32312\\
00000:0 2 3 00001:0 31212-31232 ; 00000:0 3 1 00001:0 31231-31232\\
00000:0 3 2 00001:0 31231-31232 ; 00000:0 3 3 00001:0 23121-23123\\
00000:0 4 1 00001:0 31212-31232 ; 00000:0 4 2 00001:0 21231-23231\\
00000:0 4 3 00001:0 12312-32312 ; 00000:1 0 1 00001:0 32131-32132\\
00000:1 0 2 00001:0 31321-32321 ; 00000:1 0 3 00001:0 12132-13132\\
00000:1 1 1 00001:0 13213-23213 ; 00000:1 1 2 00001:0 21321-21323\\
00000:1 1 3 00001:0 12132-32132 ; 00000:1 2 1 00001:0 13212-13232\\
00000:1 2 2 00001:0 12132-13132 ; 00000:1 2 3 00001:0 13212-13232\\
00000:1 3 1 00001:0 21321-21323 ; 00000:1 3 2 00001:0 32131-32132\\
00000:1 3 3 00001:0 21321-21323 ; 00000:1 4 1 00001:0 12132-13132\\
00000:1 4 2 00001:0 13213-23213 ; 00000:1 4 3 00001:0 13212-13232\\
00001:0 0 1 00011:0 21230-23230 ; 00001:0 0 2 00012:0 12132-13132\\
00001:0 0 3 00012:0 12132-13132 ; 00001:0 1 1 00101:0 12132-32132\\
00001:0 1 2 00102:0 21313-21323 ; 00001:0 1 3 00102:0 12313-12323\\
00001:0 2 1 00101:0 12310-32310 ; 00001:0 2 2 00102:0 12132-13132\\
00001:0 2 3 00102:0 12132-13132 ; 00001:0 3 1 00011:0 12132-13132\\
00001:0 3 2 00012:0 12132-13132 ; 00001:0 3 3 00012:0 13123-23123\\
00001:0 4 0 00000:0 31212-31232 ; 00001:0 4 0 00000:1 12132-13132\\
00001:0 4 2 00001:0 12313-12323 ; 00001:0 4 3 00001:0 31212-31232\\
00011:0 0 1 00011:0 21230-23230 ; 00011:0 0 2 00112:0 12130-12132\\
00011:0 0 3 00112:0 12120-12130 ; 00011:0 1 1 00101:0 12103-12123\\
00011:0 1 2 00121:0 13103-23103 ; 00011:0 1 3 00121:0 12102-32102\\
00011:0 2 1 00011:0 12302-32302 ; 00011:0 2 2 00112:0 12103-12123\\
00011:0 2 3 00112:0 21230-31230 ; 00011:0 3 0 00001:0 12132-13132\\
00011:0 3 2 00012:0 12102-32102 ; 00011:0 3 3 00012:0 12102-32102\\
00011:0 4 0 00001:0 12123-12323 ; 00011:0 4 2 00012:0 12103-12123\\
00011:0 4 3 00012:0 12102-32102 ; 00012:0 0 1 00121:0 21203-23203\\
00012:0 0 2 00112:0 12320-32320 ; 00012:0 0 3 00123:0 12123-13123\\
00012:0 1 1 01012:0 12103-12123 ; 00012:0 1 2 01012:0 13103-13203\\
00012:0 1 3 01023:0 12103-12123 ; 00012:0 2 1 00112:0 12320-32320\\
00012:0 2 2 00121:0 12103-12303 ; 00012:0 2 3 00123:0 12103-12123\\
00012:0 3 0 00001:0 12130-13130 ; 00012:0 3 2 00011:0 12130-13130\\
00012:0 3 3 00012:0 12120-13120 ; 00012:0 4 0 00001:0 12123-13123\\
00012:0 4 1 00011:0 12103-12123 ; 00012:0 4 3 00012:0 12120-13120\\
00101:0 0 1 00101:0 21020-23020 ; 00101:0 0 2 01012:0 12010-32010\\
00101:0 0 3 01012:0 12013-12313 ; 00101:0 1 1 00101:0 12013-12313\\
00101:0 1 2 01012:0 13010-13020 ; 00101:0 1 3 01012:0 12012-32012\\
00101:0 2 0 00001:0 12310-32310 ; 00101:0 2 2 00102:0 12012-32012\\
00101:0 2 3 00102:0 12012-32012 ; 00101:0 3 1 00011:0 12023-13023\\
00101:0 3 2 00121:0 12010-32010 ; 00101:0 3 3 00121:0 12010-32010\\
00101:0 4 0 00001:0 12032-32032 ; 00101:0 4 2 00102:0 12010-32010\\
00101:0 4 3 00102:0 12010-32010 ; 00102:0 0 1 01012:0 21013-23013\\
00102:0 0 2 00121:0 12010-12030 ; 00102:0 0 3 01023:0 12310-12320\\
00102:0 1 1 00121:0 12013-12313 ; 00102:0 1 2 01012:0 13010-23010\\
00102:0 1 3 01023:0 12010-32010 ; 00102:0 2 0 00001:0 13213-23213\\
00102:0 2 2 00101:0 12010-32010 ; 00102:0 2 3 00102:0 12010-32010\\
00102:0 3 1 00112:0 12320-32320 ; 00102:0 3 2 00112:0 12310-32310\\
00102:0 3 3 00123:0 12010-32010 ; 00102:0 4 0 00001:0 21013-23013\\
00102:0 4 1 00101:0 12010-32010 ; 00102:0 4 3 00102:0 12010-12030\\
00112:0 0 1 00102:0 21030-21230 ; 00112:0 0 2 00112:0 12023-13023\\
00112:0 0 3 00123:0 12023-13023 ; 00112:0 1 1 00012:0 12023-13023\\
00112:0 1 2 00121:0 13020-13023 ; 00112:0 1 3 00123:0 12320-32320\\
00112:0 2 0 00012:0 12320-32320 ; 00112:0 2 2 00121:0 12023-13023\\
00112:0 2 3 00123:0 13020-13023 ; 00112:0 3 0 00102:0 12320-32320\\
00112:0 3 2 00112:0 13030-13230 ; 00112:0 3 3 00123:0 13020-13023\\
00112:0 4 0 00011:0 12023-13023 ; 00112:0 4 1 00011:0 12023-13023\\
00112:0 4 3 00112:0 12320-32320 ; 00121:0 0 1 00102:0 21013-23013\\
00121:0 0 2 01012:0 12010-32010 ; 00121:0 0 3 01023:0 12013-12313\\
00121:0 1 1 00102:0 12010-32010 ; 00121:0 1 2 01012:0 13030-13230\\
00121:0 1 3 01023:0 12013-12313 ; 00121:0 2 0 00012:0 12302-12303\\
00121:0 2 2 00112:0 12010-32010 ; 00121:0 2 3 00123:0 12010-32010\\
00121:0 3 0 00101:0 12010-32010 ; 00121:0 3 1 00011:0 12032-32032\\
00121:0 3 3 00121:0 12010-32010 ; 00121:0 4 0 00012:0 12012-13012\\
00121:0 4 2 00112:0 12013-13013 ; 00121:0 4 3 00123:0 12010-32010\\
00123:0 0 1 01023:0 21010-23010 ; 00123:0 0 2 01023:0 12010-12310\\
00123:0 0 3 00123:0 12010-12310 ; 00123:0 1 1 00123:0 12012-13012\\
00123:0 1 2 01023:0 13012-13032 ; 00123:0 1 3 01023:0 12030-12032\\
00123:0 2 0 00012:0 12302-32302 ; 00123:0 2 2 00112:0 12030-32030\\
00123:0 2 3 00121:0 12010-32010 ; 00123:0 3 0 00102:0 12010-32010\\
00123:0 3 1 00112:0 12030-12032 ; 00123:0 3 3 00112:0 12012-13012\\
00123:0 4 0 00012:0 12012-13012 ; 00123:0 4 1 00121:0 12010-32010\\
00123:0 4 2 00112:0 12310-32310 ; 01012:0 0 1 00102:0 20103-23103\\
01012:0 0 2 00121:0 10103-10203 ; 01012:0 0 3 01023:0 10120-13120\\
01012:0 1 0 00012:0 12103-12123 ; 01012:0 1 2 01012:0 10120-13120\\
01012:0 1 3 01023:0 10130-10230 ; 01012:0 2 1 00012:0 10203-13203\\
01012:0 2 2 01012:0 10303-12303 ; 01012:0 2 3 01023:0 10120-13120\\
01012:0 3 0 00102:0 10123-12123 ; 01012:0 3 2 00121:0 10130-10230\\
01012:0 3 3 01023:0 10103-10203 ; 01012:0 4 0 00101:0 10103-10203\\
01012:0 4 1 00101:0 10120-13120 ; 01012:0 4 3 01012:0 10230-30230\\
01023:0 0 1 00123:0 20101-23101 ; 01023:0 0 2 01023:0 10102-30102\\
01023:0 0 3 00123:0 10130-10132 ; 01023:0 1 0 00012:0 12130-32130\\
01023:0 1 2 01012:0 10102-30102 ; 01023:0 1 3 01012:0 10102-10302\\
01023:0 2 1 00123:0 10202-13202 ; 01023:0 2 2 00123:0 10102-10302\\
01023:0 2 3 01023:0 10130-10132 ; 01023:0 3 0 00102:0 10130-30130\\
01023:0 3 1 01012:0 10102-10302 ; 01023:0 3 3 00121:0 10102-30102\\
01023:0 4 0 00102:0 10212-30212 ; 01023:0 4 1 01012:0 10102-10302\\
01023:0 4 2 00121:0 10210-30210 ; 00004:0 0 1 00014:0 21231-23231\\
00004:0 0 2 00014:0 12120-12320 ; 00004:0 0 3 00014:0 12120-12130\\
00004:0 1 1 00104:0 12120-12130 ; 00004:0 1 2 00104:0 13120-13123\\
00004:0 1 3 00104:0 12120-12130 ; 00004:0 2 1 00104:0 12310-12312\\
00004:0 2 2 00104:0 12123-13123 ; 00004:0 2 3 00104:0 12123-13123\\
00004:0 3 1 00014:0 12123-13123 ; 00004:0 3 2 00014:0 12132-13132\\
00004:0 3 3 00014:0 12123-12323 ; 00014:0 0 1 01104:0 21230-23230\\
00014:0 0 2 01124:0 12120-13120 ; 00014:0 0 3 01124:0 12102-13102\\
00014:0 1 1 01014:0 12130-32130 ; 00014:0 1 2 01024:0 13132-13232\\
00014:0 1 3 01024:0 12130-32130 ; 00014:0 2 1 00114:0 12302-12303\\
00014:0 2 2 00124:0 12120-13120 ; 00014:0 2 3 00124:0 12102-13102\\
00014:0 3 0 00004:0 12123-13123 ; 00014:0 3 2 00014:0 12130-32130\\
00014:0 3 3 00014:0 12120-13120 ; 00104:0 0 1 01014:0 21010-31010\\
00104:0 0 2 01024:0 12010-12030 ; 00104:0 0 3 01024:0 12010-12030\\
00104:0 1 1 01104:0 12010-32010 ; 00104:0 1 2 01204:0 13010-23010\\
00104:0 1 3 01204:0 12010-32010 ; 00104:0 2 0 00004:0 12310-12312\\
00104:0 2 2 00104:0 12010-32010 ; 00104:0 2 3 00104:0 12010-32010\\
00104:0 3 1 00114:0 12020-32020 ; 00104:0 3 2 00124:0 12010-12030\\
00104:0 3 3 00124:0 12012-13012 ; 00114:0 0 1 00104:0 21020-21030\\
00114:0 0 2 00124:0 12023-12323 ; 00114:0 0 3 00124:0 12020-12320\\
00114:0 1 1 00014:0 12023-13023 ; 00114:0 1 2 00124:0 13020-13023\\
00114:0 1 3 00124:0 12023-12323 ; 00114:0 2 0 00014:0 12302-12303\\
00114:0 2 2 00124:0 12020-12320 ; 00114:0 2 3 00124:0 12020-32020\\
00114:0 3 0 00104:0 12020-32020 ; 00114:0 3 2 00124:0 12030-32030\\
00114:0 3 3 00124:0 12020-12320 ; 00124:0 0 1 01024:0 21010-21310\\
00124:0 0 2 01204:0 12010-32010 ; 00124:0 0 3 01234:0 12012-13012\\
00124:0 1 1 01124:0 12010-32010 ; 00124:0 1 2 01024:0 13030-13230\\
00124:0 1 3 01234:0 12010-32010 ; 00124:0 2 0 00014:0 12302-32302\\
00124:0 2 2 00114:0 12010-12030 ; 00124:0 2 3 00124:0 12010-32010\\
00124:0 3 0 00104:0 12010-12030 ; 00124:0 3 1 00114:0 12030-32030\\
00124:0 3 3 00124:0 12012-13012 ; 01014:0 0 1 00104:0 20101-30101\\
01014:0 0 2 01024:0 10102-12102 ; 01014:0 0 3 01024:0 10103-10203\\
01014:0 1 0 00014:0 12130-32130 ; 01014:0 1 2 01024:0 10120-30120\\
01014:0 1 3 01024:0 10102-12102 ; 01014:0 2 1 00014:0 10230-30230\\
01014:0 2 2 01024:0 10102-10132 ; 01014:0 2 3 01024:0 10102-10132\\
01014:0 3 0 00104:0 10123-12123 ; 01014:0 3 2 01024:0 10103-12103\\
01014:0 3 3 01024:0 10103-12103 ; 01024:0 0 1 00124:0 20101-20301\\
01024:0 0 2 01204:0 10103-12103 ; 01024:0 0 3 01234:0 10103-12103\\
01024:0 1 0 00014:0 12102-32102 ; 01024:0 1 2 01014:0 10102-13102\\
01024:0 1 3 01024:0 10103-12103 ; 01024:0 2 1 01124:0 10203-20203\\
01024:0 2 2 00124:0 10102-13102 ; 01024:0 2 3 01234:0 10102-13102\\
01024:0 3 0 00104:0 10130-10230 ; 01024:0 3 1 01014:0 10103-12103\\
01024:0 3 3 01024:0 10102-13102 ; 01104:0 0 1 00014:0 20230-30230\\
01104:0 0 2 01124:0 10230-13230 ; 01104:0 0 3 01124:0 10213-20213\\
01104:0 1 0 00104:0 12010-32010 ; 01104:0 1 2 01204:0 10210-13210\\
01104:0 1 3 01204:0 10213-20213 ; 01104:0 2 0 00104:0 10212-10232\\
01104:0 2 2 01204:0 10310-20310 ; 01104:0 2 3 01204:0 10212-10232\\
01104:0 3 1 00014:0 10320-30320 ; 01104:0 3 2 01124:0 10212-10213\\
01104:0 3 3 01124:0 10210-13210 ; 01124:0 0 1 00014:0 20201-23201\\
01124:0 0 2 01104:0 10202-10302 ; 01124:0 0 3 01124:0 10202-10302\\
01124:0 1 0 00124:0 12010-32010 ; 01124:0 1 2 01024:0 10202-10302\\
01124:0 1 3 01234:0 10202-10302 ; 01124:0 2 0 01024:0 10203-20203\\
01124:0 2 2 00124:0 10310-12310 ; 01124:0 2 3 01234:0 10203-20203\\
01124:0 3 0 01104:0 10212-10213 ; 01124:0 3 1 00014:0 10202-10302\\
01124:0 3 3 01124:0 10203-20203 ; 01204:0 0 1 00124:0 20120-23120\\
01204:0 0 2 01024:0 10120-30120 ; 01204:0 0 3 01234:0 10120-20120\\
01204:0 1 0 00104:0 12012-13012 ; 01204:0 1 2 01104:0 10130-30130\\
01204:0 1 3 01204:0 10120-30120 ; 01204:0 2 0 00104:0 10123-12123\\
01204:0 2 1 01104:0 10310-20310 ; 01204:0 2 3 01204:0 10120-20120\\
01204:0 3 1 01024:0 10130-20130 ; 01204:0 3 2 00124:0 10130-30130\\
01204:0 3 3 01234:0 10120-30120 ; 01234:0 0 1 00124:0 20103-20303\\
01234:0 0 2 01024:0 10102-10302 ; 01234:0 0 3 01204:0 10102-10302\\
01234:0 1 0 00124:0 12012-13012 ; 01234:0 1 2 01124:0 10102-10302\\
01234:0 1 3 01024:0 10102-10302 ; 01234:0 2 0 01024:0 10103-12103\\
01234:0 2 1 01124:0 10302-30302 ; 01234:0 2 3 00124:0 10103-12103\\
01234:0 3 0 01204:0 10120-30120 ; 01234:0 3 1 01024:0 10102-10302\\
01234:0 3 2 00124:0 10102-10302\\
\end{flushleft}\endgroup
% END INTERFACES-ALL

\begin{thebibliography}{9}
\bibitem{DFFV06}
M.~Dyer, A.~D.~Flaxman, A.~M.~Frieze, and E.~Vigoda,
\newblock Randomly coloring sparse random graphs with fewer colors than the maximum degree,
\newblock \emph{Random Structures \& Algorithms} 29 (2006), no.~4, 450--465.
\newblock \doi{10.1002/rsa.20129}.

\bibitem{ChMa11}
K.~Choo and G.~MacGillivray,
\newblock Gray code numbers for graphs,
\newblock \emph{Ars Mathematica Contemporanea} 4 (2011), no.~1, 125--139.
\newblock \doi{10.26493/1855-3974.196.0df}.

\bibitem{BP16}
N.~Bousquet and G.~Perarnau,
\newblock Fast recoloring of sparse graphs,
\newblock \emph{European Journal of Combinatorics} 52A (2016), 1--11.
\newblock \doi{10.1016/j.ejc.2015.08.001}.

\bibitem{HLTH04}
H.-C.~Hsu, T.-K.~Li, J.~J.~M.~Tan, and L.-H.~Hsu,
\newblock Fault Hamiltonicity and fault Hamiltonian connectivity of the arrangement graphs,
\newblock \emph{IEEE Transactions on Computers} 53 (2004), no.~1, 39--53.
\newblock \doi{10.1109/TC.2004.1255789}.

\bibitem{BBFHMP23}
V.~Bartier, N.~Bousquet, C.~Feghali, M.~Heinrich, B.~Moore, and T.~Pierron,
\newblock Recolouring planar graphs of girth at least five,
\newblock \emph{SIAM Journal on Discrete Mathematics} 37 (2023), no.~1, 332--350.
\newblock \doi{10.1137/21M1463598}.

\bibitem{Grotzsch59}
H.~Gr\"otzsch,
\newblock Zur Theorie der diskreten Gebilde, VII: Ein Dreifarbensatz f\"ur dreikreisfreie Netze auf der Kugel,
\newblock \emph{Wiss. Z. Martin-Luther-Univ. Halle-Wittenberg, Math.-Nat. Reihe} 8 (1958/1959), 109--120.
\end{thebibliography}

\end{document}
